% !TEX program = lualatex
% ---------------------------------------------------------------------------------------------
%  Readback note 14: Auxiliary dual pairs, determinant polynomials and the carrier: the explicit matching of theta kernels
%
%  Build:  python3 tools/rb.py note <slug>
%
%  One unit of a Lean artifact, rewritten as ordinary mathematics. The prose is hand-written. Lean names,
%  line ranges and counts come from data.tex; quoted Lean statements are copied from the artifact by the build.
%
%  Rules of the text: standard terminology; every symbol is explained where it first appears and listed again
%  in the notation at the end; no letter has two meanings; a proof is complete or says what it takes for granted;
%  Lean names appear only in the LEAN lines and in the concordance.
%
%  The signature rule (the build enforces it): a function, relation or predicate that the artifact defines is
%  written only through its notation macro, which the build makes from the spec (`notation`) with one slot for
%  each explicit argument of the Lean definition. No argument is fixed and dropped, also not "for brevity" and
%  not for one section. Every argument and every hypothesis of a stated declaration stands inside its
%  statement, or in a Setting that the statement cites by number; only implicit and instance arguments (the
%  ambient objects) may stand in the conventions.
% ---------------------------------------------------------------------------------------------
\documentclass[11pt,a4paper]{article}
\ifdefined\pdfvariable\pdfvariable suppressoptionalinfo \numexpr 1+2+4+8\relax\fi
\usepackage[a4paper,hmargin=35mm,top=22mm,bottom=27mm,footskip=12mm,nohead]{geometry}

\usepackage{fontspec}
\directlua{
  luaotfload.add_fallback("gfallback", {"STIXTwoMath-Regular.otf:mode=node;", "DejaVuSans.ttf:mode=node;"})
  luaotfload.add_fallback("gmonofallback", {"DejaVuSansMono.ttf:mode=node;", "STIXTwoMath-Regular.otf:mode=node;"})
}
\setmainfont{NewCM10}[Extension=.otf,UprightFont=*-Book,ItalicFont=*-BookItalic,
  BoldFont=*-Bold,BoldItalicFont=*-BoldItalic,SmallCapsFont=*-Book,SmallCapsFeatures={Letters=SmallCaps},
  RawFeature={fallback=gfallback}]
\setsansfont{NewCMSans10}[Extension=.otf,UprightFont=*-Book,ItalicFont=*-BookOblique,
  BoldFont=*-Bold,BoldItalicFont=*-BoldOblique,RawFeature={fallback=gfallback}]
\setmonofont{JetBrainsMono}[Extension=.otf,UprightFont=*-Regular,ItalicFont=*-Italic,
  BoldFont=*-SemiBold,BoldItalicFont=*-SemiBoldItalic,Scale=MatchLowercase,
  RawFeature={-calt;-liga;fallback=gmonofallback}]
\usepackage{amsmath}
\usepackage{amsthm}
\usepackage[math-style=TeX]{unicode-math}
\setmathfont{NewCMMath-Book.otf}
\emergencystretch=2.5em\relax

\usepackage{xcolor}
\usepackage{tikz}
\usetikzlibrary{calc,arrows.meta,positioning,backgrounds}
\usepackage[protrusion=true,expansion=true]{microtype}
\usepackage{enumitem}
\usepackage{needspace}
\usepackage{booktabs}
\usepackage{array}
\usepackage{longtable}
\usepackage{fancyhdr}
\usepackage{lastpage}
\usepackage{multicol}
\usepackage[hidelinks,pdfusetitle]{hyperref}

\definecolor{ink}{HTML}{0F172A}
\definecolor{inksoft}{HTML}{475569}
\definecolor{inkfaint}{HTML}{94A3B8}
\definecolor{rulegrey}{HTML}{CBD5E1}
\definecolor{arith}{HTML}{C2410C}
\definecolor{est}{HTML}{7C3AED}
\definecolor{ana}{HTML}{0891B2}
\definecolor{ok}{HTML}{047857}
\definecolor{leancol}{HTML}{3730A3}

\title{Auxiliary dual pairs, determinant polynomials and the carrier: the explicit matching of theta kernels}
\author{grh-compact}
\date{}

% ------------------------------------------------------------------------------------- macros
\newcommand{\leanfont}{\ttfamily\fontsize{8.6}{10.5}\selectfont\color{leancol}}
\newcommand{\lean}[1]{\csname LN@#1\endcsname}
\newcommand{\fillin}[1]{{\color{arith}\itshape[#1]}}  % a place of the template that is not written yet; the build reports every one that is left
\newcommand{\leanquote}[1]{\csname LQ@#1\endcsname}
% generated by rb.py note; do not edit
\newcommand{\ConcDecls}{73}
\newcommand{\Pin}{72fead10}
\newcommand{\PinText}{that of commit {\leanfont 72fead10}}
\newcommand{\ArtifactName}{grh-compact}
\expandafter\newcommand\csname LN@FixesColsOff\endcsname{\href{https://github.com/lsernau/grh-lean/blob/72fead102f663d5c7a407964a13dc0cf9d977207/grh-compact/src/New/Rows/PartnerCut.lean\#L16-L17}{{\leanfont FixesColsOff}}}
\expandafter\newcommand\csname LN@rowStabilizer\endcsname{\href{https://github.com/lsernau/grh-lean/blob/72fead102f663d5c7a407964a13dc0cf9d977207/grh-compact/src/New/Rows/PartnerCut.lean\#L21-L23}{{\leanfont rowStabilizer}}}
\expandafter\newcommand\csname LN@columnDatum\endcsname{\href{https://github.com/lsernau/grh-lean/blob/72fead102f663d5c7a407964a13dc0cf9d977207/grh-compact/src/New/Rows/PartnerCut.lean\#L27-L28}{{\leanfont columnDatum}}}
\expandafter\newcommand\csname LN@columnDatumMul\endcsname{\href{https://github.com/lsernau/grh-lean/blob/72fead102f663d5c7a407964a13dc0cf9d977207/grh-compact/src/New/Rows/PartnerCut.lean\#L33-L43}{{\leanfont columnDatum\_mul\_ite\_eq}}}
\expandafter\newcommand\csname LN@HasWeightNormLe\endcsname{\href{https://github.com/lsernau/grh-lean/blob/72fead102f663d5c7a407964a13dc0cf9d977207/grh-compact/src/New/Rows/SemiInvariantNorm.lean\#L11-L11}{{\leanfont HasWeightNormLe}}}
\expandafter\newcommand\csname LN@HasWeightNormLeProd\endcsname{\href{https://github.com/lsernau/grh-lean/blob/72fead102f663d5c7a407964a13dc0cf9d977207/grh-compact/src/New/Rows/SemiInvariantNorm.lean\#L13-L20}{{\leanfont HasWeightNormLe.\allowbreak{}prod}}}
\expandafter\newcommand\csname LN@ChosenGL\endcsname{\href{https://github.com/lsernau/grh-lean/blob/72fead102f663d5c7a407964a13dc0cf9d977207/grh-compact/src/New/Rows/StackedDeterminants.lean\#L33-L33}{{\leanfont LevelPlaces.\allowbreak{}ChosenGL}}}
\expandafter\newcommand\csname LN@firstBlockHom\endcsname{\href{https://github.com/lsernau/grh-lean/blob/72fead102f663d5c7a407964a13dc0cf9d977207/grh-compact/src/New/Rows/StackedDeterminants.lean\#L40-L41}{{\leanfont firstBlockHom}}}
\expandafter\newcommand\csname LN@auxLevel\endcsname{\href{https://github.com/lsernau/grh-lean/blob/72fead102f663d5c7a407964a13dc0cf9d977207/grh-compact/src/New/Rows/StackedDeterminants.lean\#L45-L47}{{\leanfont auxLevel}}}
\expandafter\newcommand\csname LN@AuxPair\endcsname{\href{https://github.com/lsernau/grh-lean/blob/72fead102f663d5c7a407964a13dc0cf9d977207/grh-compact/src/New/Rows/StackedDeterminants.lean\#L58-L66}{{\leanfont AuxPair}}}
\expandafter\newcommand\csname LN@levelGroup\endcsname{\href{https://github.com/lsernau/grh-lean/blob/72fead102f663d5c7a407964a13dc0cf9d977207/grh-compact/src/New/Rows/StackedDeterminants.lean\#L76-L76}{{\leanfont AuxPair.\allowbreak{}levelGroup}}}
\expandafter\newcommand\csname LN@resDetDatum\endcsname{\href{https://github.com/lsernau/grh-lean/blob/72fead102f663d5c7a407964a13dc0cf9d977207/grh-compact/src/New/Rows/StackedDeterminants.lean\#L80-L81}{{\leanfont AuxPair.\allowbreak{}resDetDatum}}}
\expandafter\newcommand\csname LN@auxKappa\endcsname{\href{https://github.com/lsernau/grh-lean/blob/72fead102f663d5c7a407964a13dc0cf9d977207/grh-compact/src/New/Rows/StackedDeterminants.lean\#L84-L84}{{\leanfont AuxPair.\allowbreak{}κ}}}
\expandafter\newcommand\csname LN@detExp\endcsname{\href{https://github.com/lsernau/grh-lean/blob/72fead102f663d5c7a407964a13dc0cf9d977207/grh-compact/src/New/Rows/StackedDeterminants.lean\#L88-L88}{{\leanfont AuxPair.\allowbreak{}detExp}}}
\expandafter\newcommand\csname LN@auxLevelChar\endcsname{\href{https://github.com/lsernau/grh-lean/blob/72fead102f663d5c7a407964a13dc0cf9d977207/grh-compact/src/New/Rows/StackedDeterminants.lean\#L92-L94}{{\leanfont AuxPair.\allowbreak{}auxLevelChar}}}
\expandafter\newcommand\csname LN@levelChar\endcsname{\href{https://github.com/lsernau/grh-lean/blob/72fead102f663d5c7a407964a13dc0cf9d977207/grh-compact/src/New/Rows/StackedDeterminants.lean\#L98-L98}{{\leanfont AuxPair.\allowbreak{}levelChar}}}
\expandafter\newcommand\csname LN@FormV\endcsname{\href{https://github.com/lsernau/grh-lean/blob/72fead102f663d5c7a407964a13dc0cf9d977207/grh-compact/src/New/Rows/StackedDeterminants.lean\#L102-L102}{{\leanfont AuxPair.\allowbreak{}FormV}}}
\expandafter\newcommand\csname LN@chosenGLRep\endcsname{\href{https://github.com/lsernau/grh-lean/blob/72fead102f663d5c7a407964a13dc0cf9d977207/grh-compact/src/New/Rows/StackedDeterminants.lean\#L106-L111}{{\leanfont AuxPair.\allowbreak{}chosenGLRep}}}
\expandafter\newcommand\csname LN@formAction\endcsname{\href{https://github.com/lsernau/grh-lean/blob/72fead102f663d5c7a407964a13dc0cf9d977207/grh-compact/src/New/Rows/StackedDeterminants.lean\#L115-L115}{{\leanfont AuxPair.\allowbreak{}formAction}}}
\expandafter\newcommand\csname LN@tensorPoly\endcsname{\href{https://github.com/lsernau/grh-lean/blob/72fead102f663d5c7a407964a13dc0cf9d977207/grh-compact/src/New/Rows/StackedDeterminants.lean\#L119-L121}{{\leanfont AuxPair.\allowbreak{}tensorPoly}}}
\expandafter\newcommand\csname LN@tensorPolyVec\endcsname{\href{https://github.com/lsernau/grh-lean/blob/72fead102f663d5c7a407964a13dc0cf9d977207/grh-compact/src/New/Rows/StackedDeterminants.lean\#L124-L127}{{\leanfont AuxPair.\allowbreak{}tensorPoly\_vecMulVec}}}
\expandafter\newcommand\csname LN@harmonicPoly\endcsname{\href{https://github.com/lsernau/grh-lean/blob/72fead102f663d5c7a407964a13dc0cf9d977207/grh-compact/src/New/Rows/StackedDeterminants.lean\#L131-L131}{{\leanfont AuxPair.\allowbreak{}harmonicPoly}}}
\expandafter\newcommand\csname LN@auxA\endcsname{\href{https://github.com/lsernau/grh-lean/blob/72fead102f663d5c7a407964a13dc0cf9d977207/grh-compact/src/New/Rows/StackedDeterminants.lean\#L134-L134}{{\leanfont AuxPair.\allowbreak{}a}}}
\expandafter\newcommand\csname LN@auxB\endcsname{\href{https://github.com/lsernau/grh-lean/blob/72fead102f663d5c7a407964a13dc0cf9d977207/grh-compact/src/New/Rows/StackedDeterminants.lean\#L138-L138}{{\leanfont AuxPair.\allowbreak{}b}}}
\expandafter\newcommand\csname LN@pairThetaKernel\endcsname{\href{https://github.com/lsernau/grh-lean/blob/72fead102f663d5c7a407964a13dc0cf9d977207/grh-compact/src/New/Rows/StackedDeterminants.lean\#L142-L142}{{\leanfont AuxPair.\allowbreak{}pairThetaKernel}}}
\expandafter\newcommand\csname LN@auxDatum\endcsname{\href{https://github.com/lsernau/grh-lean/blob/72fead102f663d5c7a407964a13dc0cf9d977207/grh-compact/src/New/Rows/StackedDeterminants.lean\#L146-L147}{{\leanfont AuxPair.\allowbreak{}auxDatum}}}
\expandafter\newcommand\csname LN@fockDatum\endcsname{\href{https://github.com/lsernau/grh-lean/blob/72fead102f663d5c7a407964a13dc0cf9d977207/grh-compact/src/New/Rows/StackedDeterminants.lean\#L151-L151}{{\leanfont AuxPair.\allowbreak{}fockDatum}}}
\expandafter\newcommand\csname LN@conjKernel\endcsname{\href{https://github.com/lsernau/grh-lean/blob/72fead102f663d5c7a407964a13dc0cf9d977207/grh-compact/src/New/Rows/StackedDeterminants.lean\#L155-L156}{{\leanfont AuxPair.\allowbreak{}conjKernel}}}
\expandafter\newcommand\csname LN@descent\endcsname{\href{https://github.com/lsernau/grh-lean/blob/72fead102f663d5c7a407964a13dc0cf9d977207/grh-compact/src/New/Rows/StackedDeterminants.lean\#L160-L161}{{\leanfont AuxPair.\allowbreak{}θ}}}
\expandafter\newcommand\csname LN@heckePolyOnForms\endcsname{\href{https://github.com/lsernau/grh-lean/blob/72fead102f663d5c7a407964a13dc0cf9d977207/grh-compact/src/New/Rows/StackedDeterminants.lean\#L167-L170}{{\leanfont AuxPair.\allowbreak{}heckePolyOnForms}}}
\expandafter\newcommand\csname LN@heckePolySmul\endcsname{\href{https://github.com/lsernau/grh-lean/blob/72fead102f663d5c7a407964a13dc0cf9d977207/grh-compact/src/New/Rows/StackedDeterminants.lean\#L173-L175}{{\leanfont AuxPair.\allowbreak{}heckePolyOnForms\_eq\_smul}}}
\expandafter\newcommand\csname LN@rankDiff\endcsname{\href{https://github.com/lsernau/grh-lean/blob/72fead102f663d5c7a407964a13dc0cf9d977207/grh-compact/src/New/Rows/StackedDeterminants.lean\#L178-L178}{{\leanfont AuxPair.\allowbreak{}rankDiff}}}
\expandafter\newcommand\csname LN@testValue\endcsname{\href{https://github.com/lsernau/grh-lean/blob/72fead102f663d5c7a407964a13dc0cf9d977207/grh-compact/src/New/Rows/StackedDeterminants.lean\#L182-L182}{{\leanfont AuxPair.\allowbreak{}testValue}}}
\expandafter\newcommand\csname LN@Lifts\endcsname{\href{https://github.com/lsernau/grh-lean/blob/72fead102f663d5c7a407964a13dc0cf9d977207/grh-compact/src/New/Rows/StackedDeterminants.lean\#L186-L188}{{\leanfont AuxPair.\allowbreak{}Lifts}}}
\expandafter\newcommand\csname LN@chosenModulusSq\endcsname{\href{https://github.com/lsernau/grh-lean/blob/72fead102f663d5c7a407964a13dc0cf9d977207/grh-compact/src/New/Rows/StackedDeterminants.lean\#L199-L199}{{\leanfont LevelPlaces.\allowbreak{}chosenModulusSq}}}
\expandafter\newcommand\csname LN@mapBlock\endcsname{\href{https://github.com/lsernau/grh-lean/blob/72fead102f663d5c7a407964a13dc0cf9d977207/grh-compact/src/New/Rows/StackedDeterminants.lean\#L201-L203}{{\leanfont map\_block\_eq\_blockGL}}}
\expandafter\newcommand\csname LN@levelModulusIdealLe\endcsname{\href{https://github.com/lsernau/grh-lean/blob/72fead102f663d5c7a407964a13dc0cf9d977207/grh-compact/src/New/Rows/StackedDeterminants.lean\#L205-L206}{{\leanfont levelModulus\_ideal\_le}}}
\expandafter\newcommand\csname LN@placesLevelModulus\endcsname{\href{https://github.com/lsernau/grh-lean/blob/72fead102f663d5c7a407964a13dc0cf9d977207/grh-compact/src/New/Rows/StackedDeterminants.lean\#L208-L210}{{\leanfont places\_levelModulus\_false}}}
\expandafter\newcommand\csname LN@reduceToRzero\endcsname{\href{https://github.com/lsernau/grh-lean/blob/72fead102f663d5c7a407964a13dc0cf9d977207/grh-compact/src/New/Rows/StackedDeterminants.lean\#L213-L214}{{\leanfont reduceToR₀}}}
\expandafter\newcommand\csname LN@liftModSupported\endcsname{\href{https://github.com/lsernau/grh-lean/blob/72fead102f663d5c7a407964a13dc0cf9d977207/grh-compact/src/New/Rows/StackedDeterminants.lean\#L216-L217}{{\leanfont liftMod\_mem\_supportedOn}}}
\expandafter\newcommand\csname LN@atPlaceTwist\endcsname{\href{https://github.com/lsernau/grh-lean/blob/72fead102f663d5c7a407964a13dc0cf9d977207/grh-compact/src/New/Rows/StackedDeterminants.lean\#L224-L225}{{\leanfont atPlace\_twistPlace\_eq\_one}}}
\expandafter\newcommand\csname LN@DividedSplitting\endcsname{\href{https://github.com/lsernau/grh-lean/blob/72fead102f663d5c7a407964a13dc0cf9d977207/grh-compact/src/New/Rows/StackedDeterminants.lean\#L229-L232}{{\leanfont DividedSplitting}}}
\expandafter\newcommand\csname LN@genIdele\endcsname{\href{https://github.com/lsernau/grh-lean/blob/72fead102f663d5c7a407964a13dc0cf9d977207/grh-compact/src/New/Rows/StackedDeterminants.lean\#L236-L236}{{\leanfont AuxPair.\allowbreak{}genIdele}}}
\expandafter\newcommand\csname LN@genIntAdele\endcsname{\href{https://github.com/lsernau/grh-lean/blob/72fead102f663d5c7a407964a13dc0cf9d977207/grh-compact/src/New/Rows/StackedDeterminants.lean\#L240-L240}{{\leanfont AuxPair.\allowbreak{}genIntAdele}}}
\expandafter\newcommand\csname LN@genIdeleApply\endcsname{\href{https://github.com/lsernau/grh-lean/blob/72fead102f663d5c7a407964a13dc0cf9d977207/grh-compact/src/New/Rows/StackedDeterminants.lean\#L244-L245}{{\leanfont genIdele\_apply\_of\_notMem}}}
\expandafter\newcommand\csname LN@reductionZero\endcsname{\href{https://github.com/lsernau/grh-lean/blob/72fead102f663d5c7a407964a13dc0cf9d977207/grh-compact/src/New/Rows/StackedDeterminants.lean\#L247-L250}{{\leanfont reduction\_levelModulus\_eq\_zero\_iff}}}
\expandafter\newcommand\csname LN@genResidue\endcsname{\href{https://github.com/lsernau/grh-lean/blob/72fead102f663d5c7a407964a13dc0cf9d977207/grh-compact/src/New/Rows/StackedDeterminants.lean\#L256-L256}{{\leanfont DividedSplitting.\allowbreak{}genResidue}}}
\expandafter\newcommand\csname LN@genResidueSmul\endcsname{\href{https://github.com/lsernau/grh-lean/blob/72fead102f663d5c7a407964a13dc0cf9d977207/grh-compact/src/New/Rows/StackedDeterminants.lean\#L258-L266}{{\leanfont DividedSplitting.\allowbreak{}genResidue\_smul\_eq\_iff}}}
\expandafter\newcommand\csname LN@stackedDet\endcsname{\href{https://github.com/lsernau/grh-lean/blob/72fead102f663d5c7a407964a13dc0cf9d977207/grh-compact/src/New/Rows/StackedDeterminants.lean\#L277-L277}{{\leanfont stackedDet}}}
\expandafter\newcommand\csname LN@unitModulus\endcsname{\href{https://github.com/lsernau/grh-lean/blob/72fead102f663d5c7a407964a13dc0cf9d977207/grh-compact/src/New/Rows/StackedDeterminants.lean\#L280-L280}{{\leanfont unitModulus}}}
\expandafter\newcommand\csname LN@carrierForm\endcsname{\href{https://github.com/lsernau/grh-lean/blob/72fead102f663d5c7a407964a13dc0cf9d977207/grh-compact/src/New/Rows/StackedDeterminants.lean\#L289-L291}{{\leanfont carrierForm}}}
\expandafter\newcommand\csname LN@carrierXi\endcsname{\href{https://github.com/lsernau/grh-lean/blob/72fead102f663d5c7a407964a13dc0cf9d977207/grh-compact/src/New/Rows/StackedDeterminants.lean\#L297-L298}{{\leanfont carrierXi}}}
\expandafter\newcommand\csname LN@detFromCols\endcsname{\href{https://github.com/lsernau/grh-lean/blob/72fead102f663d5c7a407964a13dc0cf9d977207/grh-compact/src/New/Rows/StackedDeterminants.lean\#L304-L313}{{\leanfont exists\_det\_fromCols}}}
\expandafter\newcommand\csname LN@twoSidedSubst\endcsname{\href{https://github.com/lsernau/grh-lean/blob/72fead102f663d5c7a407964a13dc0cf9d977207/grh-compact/src/New/Rows/StackedDeterminants.lean\#L317-L318}{{\leanfont twoSidedSubst}}}
\expandafter\newcommand\csname LN@inFirstCols\endcsname{\href{https://github.com/lsernau/grh-lean/blob/72fead102f663d5c7a407964a13dc0cf9d977207/grh-compact/src/New/Rows/StackedDeterminants.lean\#L322-L323}{{\leanfont inFirstCols}}}
\expandafter\newcommand\csname LN@stackedDetProd\endcsname{\href{https://github.com/lsernau/grh-lean/blob/72fead102f663d5c7a407964a13dc0cf9d977207/grh-compact/src/New/Rows/StackedDeterminants.lean\#L327-L327}{{\leanfont stackedDetProd}}}
\expandafter\newcommand\csname LN@substStackedDetProd\endcsname{\href{https://github.com/lsernau/grh-lean/blob/72fead102f663d5c7a407964a13dc0cf9d977207/grh-compact/src/New/Rows/StackedDeterminants.lean\#L332-L337}{{\leanfont twoSidedSubst\_stackedDetProd}}}
\expandafter\newcommand\csname LN@firstColsCoord\endcsname{\href{https://github.com/lsernau/grh-lean/blob/72fead102f663d5c7a407964a13dc0cf9d977207/grh-compact/src/New/Rows/StackedDeterminants.lean\#L342-L343}{{\leanfont firstColsCoord}}}
\expandafter\newcommand\csname LN@substFirstColsCoord\endcsname{\href{https://github.com/lsernau/grh-lean/blob/72fead102f663d5c7a407964a13dc0cf9d977207/grh-compact/src/New/Rows/StackedDeterminants.lean\#L345-L353}{{\leanfont twoSidedSubst\_firstColsCoord}}}
\expandafter\newcommand\csname LN@blockSigns\endcsname{\href{https://github.com/lsernau/grh-lean/blob/72fead102f663d5c7a407964a13dc0cf9d977207/grh-compact/src/New/Rows/RealMatching.lean\#L20-L20}{{\leanfont blockSigns}}}
\expandafter\newcommand\csname LN@posBlockEmb\endcsname{\href{https://github.com/lsernau/grh-lean/blob/72fead102f663d5c7a407964a13dc0cf9d977207/grh-compact/src/New/Rows/RealMatching.lean\#L24-L25}{{\leanfont posBlockEmb}}}
\expandafter\newcommand\csname LN@posBlockEmbSpec\endcsname{\href{https://github.com/lsernau/grh-lean/blob/72fead102f663d5c7a407964a13dc0cf9d977207/grh-compact/src/New/Rows/RealMatching.lean\#L28-L31}{{\leanfont posBlockEmb\_spec}}}
\expandafter\newcommand\csname LN@cardNegBlockSigns\endcsname{\href{https://github.com/lsernau/grh-lean/blob/72fead102f663d5c7a407964a13dc0cf9d977207/grh-compact/src/New/Rows/RealMatching.lean\#L33-L34}{{\leanfont card\_neg\_blockSigns}}}
\expandafter\newcommand\csname LN@matchedPoly\endcsname{\href{https://github.com/lsernau/grh-lean/blob/72fead102f663d5c7a407964a13dc0cf9d977207/grh-compact/src/New/Rows/RealMatching.lean\#L38-L39}{{\leanfont matchedPoly}}}
\expandafter\newcommand\csname LN@posFockVector\endcsname{\href{https://github.com/lsernau/grh-lean/blob/72fead102f663d5c7a407964a13dc0cf9d977207/grh-compact/src/New/Rows/RealMatching.lean\#L43-L44}{{\leanfont posFockVector}}}
\expandafter\newcommand\csname LN@realWeilStackedDetProd\endcsname{\href{https://github.com/lsernau/grh-lean/blob/72fead102f663d5c7a407964a13dc0cf9d977207/grh-compact/src/New/Rows/RealMatching.lean\#L52-L56}{{\leanfont realPairWeil\_stackedDetProd}}}
\expandafter\newcommand\csname LN@realmatching\endcsname{\href{https://github.com/lsernau/grh-lean/blob/72fead102f663d5c7a407964a13dc0cf9d977207/grh-compact/src/New/Rows/RealMatching.lean\#L61-L113}{{\leanfont real\_matching}}}
\expandafter\newcommand\csname LQ@realmatching\endcsname{\leanquoteline{0}{lemma real\_matching : ∃ Q₀, ∀ v, ∃ κ, ∀ (ℓ : \_ → \_ →L[ℂ] ℂ) (ℓ' : \_ →ₗ[ℂ] ℂ), (∀ g x w,}\leanquoteline{2}{ℓ g (realPairWeil γ 0 x w) = ℓ (g * x.1) w) → (∀ x Q, realPairWeil (s := blockSigns f) γ 0 (1,}\leanquoteline{2}{x) (polyVector \_ Q₀) = polyVector \_ Q → ℓ x v = ℓ' Q) → ∀ η,}\leanquoteline{2}{κ * ℓ 1 (posFockVector j (stackedDetProd e η)) = ℓ' (matchedPoly (ι := ι) f j η)}}
\expandafter\newcommand\csname LN@isRapidModKernel\endcsname{\href{https://github.com/lsernau/grh-lean/blob/72fead102f663d5c7a407964a13dc0cf9d977207/grh-compact/src/New/Rows/Carrier.lean\#L42-L45}{{\leanfont isRapid\_modKernel}}}
\expandafter\newcommand\csname LN@explicitMatchingOfSum\endcsname{\href{https://github.com/lsernau/grh-lean/blob/72fead102f663d5c7a407964a13dc0cf9d977207/grh-compact/src/New/Rows/Carrier.lean\#L52-L75}{{\leanfont explicit\_matching\_of\_sum}}}
\expandafter\newcommand\csname LN@classSumRowStabilizer\endcsname{\href{https://github.com/lsernau/grh-lean/blob/72fead102f663d5c7a407964a13dc0cf9d977207/grh-compact/src/New/Rows/Carrier.lean\#L80-L103}{{\leanfont stdClassSum\_rowStabilizer\_eq}}}
\expandafter\newcommand\csname LN@explicitmatching\endcsname{\href{https://github.com/lsernau/grh-lean/blob/72fead102f663d5c7a407964a13dc0cf9d977207/grh-compact/src/New/Rows/Carrier.lean\#L112-L146}{{\leanfont explicit\_matching}}}
\expandafter\newcommand\csname LQ@explicitmatching\endcsname{\leanquoteline{0}{lemma explicit\_matching \{E\} [Field E] [CMType E] [WeilData E] \{n m a b p l\} [Fintype n] [Fintype m]}\leanquoteline{2}{[Fintype a] [Fintype b] [Fintype p] [Fintype l] \{s\} (W : Space E n s) (f : p ⊕ l ≃ b) (V' : Space}\leanquoteline{2}{E \_ (blockSigns f)) σ e (j : m ⊕ l ≃ n) (hj : s = blockSigns j) \{d\} (hd : card a = card b + 2 * d)}\leanquoteline{2}{(hn : 2 * card l + 2 ≤ card n) ι [∀ τ, Fintype (ι τ)] : ∃ C, ∀ \{𝔪 α γ χ\}, IsTameFor 𝔪 α γ χ →}\leanquoteline{2}{(∀ τ, γ τ = card l - card (ι τ)) → ∃ c, ‖c‖ \textasciicircum{} 2 = C / lValueProd 𝔪.places}\leanquoteline{2}{(α⁻¹ * upsilonU1 \textasciicircum{} (card b + d)) d ∧ ∀ η \{Φ Ψ\}, IsPartner ⊤ σ χ⁻¹ Φ Ψ → ∀ β g, c * stdClassSum}\leanquoteline{2}{(Kf E a) (modKernelI W α β 𝔪 Φ (renameRowsCols id (j ∘ .inl) (stackedDetProd e η)) g ·) = ∫ x,}\leanquoteline{2}{modKernel V' W (α / β) α 𝔪 Ψ (matchedPoly (ι := ι) f j η) x g ∂V'.quotMeasure}}
\expandafter\newcommand\csname LN@classSumDet\endcsname{\href{https://github.com/lsernau/grh-lean/blob/72fead102f663d5c7a407964a13dc0cf9d977207/grh-compact/src/New/Rows/Carrier.lean\#L159-L171}{{\leanfont stdClassSum\_det\_modKernel\_eq}}}
\expandafter\newcommand\csname LQ@classSumDet\endcsname{\leanquoteline{0}{lemma stdClassSum\_det\_modKernel\_eq (he : Even m) (c : \_ → ℕ) : ∃ C, ∀ \{𝔪 α χ\},}\leanquoteline{2}{IsTameFor 𝔪 α (-c ·) χ → ∃ γ, ‖γ‖ \textasciicircum{} 2 = C / lValueProd}\leanquoteline{2}{𝔪.places (α⁻¹ * upsilonU1 \textasciicircum{} (m / 2)) (m / 2) ∧ ∀ g, γ * stdClassSum (Kf E (Fin m))}\leanquoteline{2}{(modKernelI (stdSpace ..) α α 𝔪 (χ ·.det) (∏ τ, letterMinor τ id id \textasciicircum{} c τ) g) = 1}}
\expandafter\newcommand\csname LN@carrierIdentity\endcsname{\href{https://github.com/lsernau/grh-lean/blob/72fead102f663d5c7a407964a13dc0cf9d977207/grh-compact/src/New/Rows/Carrier.lean\#L178-L199}{{\leanfont stdClassSum\_carrierXi\_eq\_carrierForm}}}
\expandafter\newcommand\csname LQ@carrierIdentity\endcsname{\leanquoteline{0}{lemma stdClassSum\_carrierXi\_eq\_carrierForm (he : Even m) : ∃ C, ∀ \{𝔪 c χ\}, IsTameFor 𝔪 c 0 χ →}\leanquoteline{2}{∃ γ, ‖γ‖ \textasciicircum{} 2 = C / (lValueProd 𝔪.places (c⁻¹ * upsilonU1 \textasciicircum{} m) m * lValueProd 𝔪.places}\leanquoteline{2}{(c * upsilonU1 \textasciicircum{} (m / 2)) (m / 2)) ∧ ∀ x, γ * stdClassSum (Kf ..) (fun h ↦ carrierXi m c h *}\leanquoteline{2}{modKernelI (signedSpace ..) (upsilonU1 \textasciicircum{} m) c 𝔪 (χ ·.det) 1 x h) = carrierForm E m x}}
\newcommand{\ConcordanceRows}{\crow{FixesColsOff}{\href{https://github.com/lsernau/grh-lean/blob/72fead102f663d5c7a407964a13dc0cf9d977207/grh-compact/src/New/Rows/PartnerCut.lean\#L16-L17}{{\leanfont FixesColsOff}}}{PartnerCut.lean}{16--17}
\crow{rowStabilizer}{\href{https://github.com/lsernau/grh-lean/blob/72fead102f663d5c7a407964a13dc0cf9d977207/grh-compact/src/New/Rows/PartnerCut.lean\#L21-L23}{{\leanfont rowStabilizer}}}{PartnerCut.lean}{21--23}
\crow{columnDatum}{\href{https://github.com/lsernau/grh-lean/blob/72fead102f663d5c7a407964a13dc0cf9d977207/grh-compact/src/New/Rows/PartnerCut.lean\#L27-L28}{{\leanfont columnDatum}}}{PartnerCut.lean}{27--28}
\crow{columnDatumMul}{\href{https://github.com/lsernau/grh-lean/blob/72fead102f663d5c7a407964a13dc0cf9d977207/grh-compact/src/New/Rows/PartnerCut.lean\#L33-L43}{{\leanfont columnDatum\_mul\_ite\_eq}}}{PartnerCut.lean}{33--43}
\crow{HasWeightNormLe}{\href{https://github.com/lsernau/grh-lean/blob/72fead102f663d5c7a407964a13dc0cf9d977207/grh-compact/src/New/Rows/SemiInvariantNorm.lean\#L11-L11}{{\leanfont HasWeightNormLe}}}{SemiInvariantNorm.lean}{11--11}
\crow{HasWeightNormLeProd}{\href{https://github.com/lsernau/grh-lean/blob/72fead102f663d5c7a407964a13dc0cf9d977207/grh-compact/src/New/Rows/SemiInvariantNorm.lean\#L13-L20}{{\leanfont HasWeightNormLe.\allowbreak{}prod}}}{SemiInvariantNorm.lean}{13--20}
\crow{ChosenGL}{\href{https://github.com/lsernau/grh-lean/blob/72fead102f663d5c7a407964a13dc0cf9d977207/grh-compact/src/New/Rows/StackedDeterminants.lean\#L33-L33}{{\leanfont LevelPlaces.\allowbreak{}ChosenGL}}}{StackedDeterminants.lean}{33--33}
\crow{firstBlockHom}{\href{https://github.com/lsernau/grh-lean/blob/72fead102f663d5c7a407964a13dc0cf9d977207/grh-compact/src/New/Rows/StackedDeterminants.lean\#L40-L41}{{\leanfont firstBlockHom}}}{StackedDeterminants.lean}{40--41}
\crow{auxLevel}{\href{https://github.com/lsernau/grh-lean/blob/72fead102f663d5c7a407964a13dc0cf9d977207/grh-compact/src/New/Rows/StackedDeterminants.lean\#L45-L47}{{\leanfont auxLevel}}}{StackedDeterminants.lean}{45--47}
\crow{AuxPair}{\href{https://github.com/lsernau/grh-lean/blob/72fead102f663d5c7a407964a13dc0cf9d977207/grh-compact/src/New/Rows/StackedDeterminants.lean\#L58-L66}{{\leanfont AuxPair}}}{StackedDeterminants.lean}{58--66}
\crow{levelGroup}{\href{https://github.com/lsernau/grh-lean/blob/72fead102f663d5c7a407964a13dc0cf9d977207/grh-compact/src/New/Rows/StackedDeterminants.lean\#L76-L76}{{\leanfont AuxPair.\allowbreak{}levelGroup}}}{StackedDeterminants.lean}{76--76}
\crow{resDetDatum}{\href{https://github.com/lsernau/grh-lean/blob/72fead102f663d5c7a407964a13dc0cf9d977207/grh-compact/src/New/Rows/StackedDeterminants.lean\#L80-L81}{{\leanfont AuxPair.\allowbreak{}resDetDatum}}}{StackedDeterminants.lean}{80--81}
\crow{auxKappa}{\href{https://github.com/lsernau/grh-lean/blob/72fead102f663d5c7a407964a13dc0cf9d977207/grh-compact/src/New/Rows/StackedDeterminants.lean\#L84-L84}{{\leanfont AuxPair.\allowbreak{}κ}}}{StackedDeterminants.lean}{84--84}
\crow{detExp}{\href{https://github.com/lsernau/grh-lean/blob/72fead102f663d5c7a407964a13dc0cf9d977207/grh-compact/src/New/Rows/StackedDeterminants.lean\#L88-L88}{{\leanfont AuxPair.\allowbreak{}detExp}}}{StackedDeterminants.lean}{88--88}
\crow{auxLevelChar}{\href{https://github.com/lsernau/grh-lean/blob/72fead102f663d5c7a407964a13dc0cf9d977207/grh-compact/src/New/Rows/StackedDeterminants.lean\#L92-L94}{{\leanfont AuxPair.\allowbreak{}auxLevelChar}}}{StackedDeterminants.lean}{92--94}
\crow{levelChar}{\href{https://github.com/lsernau/grh-lean/blob/72fead102f663d5c7a407964a13dc0cf9d977207/grh-compact/src/New/Rows/StackedDeterminants.lean\#L98-L98}{{\leanfont AuxPair.\allowbreak{}levelChar}}}{StackedDeterminants.lean}{98--98}
\crow{FormV}{\href{https://github.com/lsernau/grh-lean/blob/72fead102f663d5c7a407964a13dc0cf9d977207/grh-compact/src/New/Rows/StackedDeterminants.lean\#L102-L102}{{\leanfont AuxPair.\allowbreak{}FormV}}}{StackedDeterminants.lean}{102--102}
\crow{chosenGLRep}{\href{https://github.com/lsernau/grh-lean/blob/72fead102f663d5c7a407964a13dc0cf9d977207/grh-compact/src/New/Rows/StackedDeterminants.lean\#L106-L111}{{\leanfont AuxPair.\allowbreak{}chosenGLRep}}}{StackedDeterminants.lean}{106--111}
\crow{formAction}{\href{https://github.com/lsernau/grh-lean/blob/72fead102f663d5c7a407964a13dc0cf9d977207/grh-compact/src/New/Rows/StackedDeterminants.lean\#L115-L115}{{\leanfont AuxPair.\allowbreak{}formAction}}}{StackedDeterminants.lean}{115--115}
\crow{tensorPoly}{\href{https://github.com/lsernau/grh-lean/blob/72fead102f663d5c7a407964a13dc0cf9d977207/grh-compact/src/New/Rows/StackedDeterminants.lean\#L119-L121}{{\leanfont AuxPair.\allowbreak{}tensorPoly}}}{StackedDeterminants.lean}{119--121}
\crow{tensorPolyVec}{\href{https://github.com/lsernau/grh-lean/blob/72fead102f663d5c7a407964a13dc0cf9d977207/grh-compact/src/New/Rows/StackedDeterminants.lean\#L124-L127}{{\leanfont AuxPair.\allowbreak{}tensorPoly\_vecMulVec}}}{StackedDeterminants.lean}{124--127}
\crow{harmonicPoly}{\href{https://github.com/lsernau/grh-lean/blob/72fead102f663d5c7a407964a13dc0cf9d977207/grh-compact/src/New/Rows/StackedDeterminants.lean\#L131-L131}{{\leanfont AuxPair.\allowbreak{}harmonicPoly}}}{StackedDeterminants.lean}{131--131}
\crow{auxA}{\href{https://github.com/lsernau/grh-lean/blob/72fead102f663d5c7a407964a13dc0cf9d977207/grh-compact/src/New/Rows/StackedDeterminants.lean\#L134-L134}{{\leanfont AuxPair.\allowbreak{}a}}}{StackedDeterminants.lean}{134--134}
\crow{auxB}{\href{https://github.com/lsernau/grh-lean/blob/72fead102f663d5c7a407964a13dc0cf9d977207/grh-compact/src/New/Rows/StackedDeterminants.lean\#L138-L138}{{\leanfont AuxPair.\allowbreak{}b}}}{StackedDeterminants.lean}{138--138}
\crow{pairThetaKernel}{\href{https://github.com/lsernau/grh-lean/blob/72fead102f663d5c7a407964a13dc0cf9d977207/grh-compact/src/New/Rows/StackedDeterminants.lean\#L142-L142}{{\leanfont AuxPair.\allowbreak{}pairThetaKernel}}}{StackedDeterminants.lean}{142--142}
\crow{auxDatum}{\href{https://github.com/lsernau/grh-lean/blob/72fead102f663d5c7a407964a13dc0cf9d977207/grh-compact/src/New/Rows/StackedDeterminants.lean\#L146-L147}{{\leanfont AuxPair.\allowbreak{}auxDatum}}}{StackedDeterminants.lean}{146--147}
\crow{fockDatum}{\href{https://github.com/lsernau/grh-lean/blob/72fead102f663d5c7a407964a13dc0cf9d977207/grh-compact/src/New/Rows/StackedDeterminants.lean\#L151-L151}{{\leanfont AuxPair.\allowbreak{}fockDatum}}}{StackedDeterminants.lean}{151--151}
\crow{conjKernel}{\href{https://github.com/lsernau/grh-lean/blob/72fead102f663d5c7a407964a13dc0cf9d977207/grh-compact/src/New/Rows/StackedDeterminants.lean\#L155-L156}{{\leanfont AuxPair.\allowbreak{}conjKernel}}}{StackedDeterminants.lean}{155--156}
\crow{descent}{\href{https://github.com/lsernau/grh-lean/blob/72fead102f663d5c7a407964a13dc0cf9d977207/grh-compact/src/New/Rows/StackedDeterminants.lean\#L160-L161}{{\leanfont AuxPair.\allowbreak{}θ}}}{StackedDeterminants.lean}{160--161}
\crow{heckePolyOnForms}{\href{https://github.com/lsernau/grh-lean/blob/72fead102f663d5c7a407964a13dc0cf9d977207/grh-compact/src/New/Rows/StackedDeterminants.lean\#L167-L170}{{\leanfont AuxPair.\allowbreak{}heckePolyOnForms}}}{StackedDeterminants.lean}{167--170}
\crow{heckePolySmul}{\href{https://github.com/lsernau/grh-lean/blob/72fead102f663d5c7a407964a13dc0cf9d977207/grh-compact/src/New/Rows/StackedDeterminants.lean\#L173-L175}{{\leanfont AuxPair.\allowbreak{}heckePolyOnForms\_eq\_smul}}}{StackedDeterminants.lean}{173--175}
\crow{rankDiff}{\href{https://github.com/lsernau/grh-lean/blob/72fead102f663d5c7a407964a13dc0cf9d977207/grh-compact/src/New/Rows/StackedDeterminants.lean\#L178-L178}{{\leanfont AuxPair.\allowbreak{}rankDiff}}}{StackedDeterminants.lean}{178--178}
\crow{testValue}{\href{https://github.com/lsernau/grh-lean/blob/72fead102f663d5c7a407964a13dc0cf9d977207/grh-compact/src/New/Rows/StackedDeterminants.lean\#L182-L182}{{\leanfont AuxPair.\allowbreak{}testValue}}}{StackedDeterminants.lean}{182--182}
\crow{Lifts}{\href{https://github.com/lsernau/grh-lean/blob/72fead102f663d5c7a407964a13dc0cf9d977207/grh-compact/src/New/Rows/StackedDeterminants.lean\#L186-L188}{{\leanfont AuxPair.\allowbreak{}Lifts}}}{StackedDeterminants.lean}{186--188}
\crow{chosenModulusSq}{\href{https://github.com/lsernau/grh-lean/blob/72fead102f663d5c7a407964a13dc0cf9d977207/grh-compact/src/New/Rows/StackedDeterminants.lean\#L199-L199}{{\leanfont LevelPlaces.\allowbreak{}chosenModulusSq}}}{StackedDeterminants.lean}{199--199}
\crow{mapBlock}{\href{https://github.com/lsernau/grh-lean/blob/72fead102f663d5c7a407964a13dc0cf9d977207/grh-compact/src/New/Rows/StackedDeterminants.lean\#L201-L203}{{\leanfont map\_block\_eq\_blockGL}}}{StackedDeterminants.lean}{201--203}
\crow{levelModulusIdealLe}{\href{https://github.com/lsernau/grh-lean/blob/72fead102f663d5c7a407964a13dc0cf9d977207/grh-compact/src/New/Rows/StackedDeterminants.lean\#L205-L206}{{\leanfont levelModulus\_ideal\_le}}}{StackedDeterminants.lean}{205--206}
\crow{placesLevelModulus}{\href{https://github.com/lsernau/grh-lean/blob/72fead102f663d5c7a407964a13dc0cf9d977207/grh-compact/src/New/Rows/StackedDeterminants.lean\#L208-L210}{{\leanfont places\_levelModulus\_false}}}{StackedDeterminants.lean}{208--210}
\crow{reduceToRzero}{\href{https://github.com/lsernau/grh-lean/blob/72fead102f663d5c7a407964a13dc0cf9d977207/grh-compact/src/New/Rows/StackedDeterminants.lean\#L213-L214}{{\leanfont reduceToR₀}}}{StackedDeterminants.lean}{213--214}
\crow{liftModSupported}{\href{https://github.com/lsernau/grh-lean/blob/72fead102f663d5c7a407964a13dc0cf9d977207/grh-compact/src/New/Rows/StackedDeterminants.lean\#L216-L217}{{\leanfont liftMod\_mem\_supportedOn}}}{StackedDeterminants.lean}{216--217}
\crow{atPlaceTwist}{\href{https://github.com/lsernau/grh-lean/blob/72fead102f663d5c7a407964a13dc0cf9d977207/grh-compact/src/New/Rows/StackedDeterminants.lean\#L224-L225}{{\leanfont atPlace\_twistPlace\_eq\_one}}}{StackedDeterminants.lean}{224--225}
\crow{DividedSplitting}{\href{https://github.com/lsernau/grh-lean/blob/72fead102f663d5c7a407964a13dc0cf9d977207/grh-compact/src/New/Rows/StackedDeterminants.lean\#L229-L232}{{\leanfont DividedSplitting}}}{StackedDeterminants.lean}{229--232}
\crow{genIdele}{\href{https://github.com/lsernau/grh-lean/blob/72fead102f663d5c7a407964a13dc0cf9d977207/grh-compact/src/New/Rows/StackedDeterminants.lean\#L236-L236}{{\leanfont AuxPair.\allowbreak{}genIdele}}}{StackedDeterminants.lean}{236--236}
\crow{genIntAdele}{\href{https://github.com/lsernau/grh-lean/blob/72fead102f663d5c7a407964a13dc0cf9d977207/grh-compact/src/New/Rows/StackedDeterminants.lean\#L240-L240}{{\leanfont AuxPair.\allowbreak{}genIntAdele}}}{StackedDeterminants.lean}{240--240}
\crow{genIdeleApply}{\href{https://github.com/lsernau/grh-lean/blob/72fead102f663d5c7a407964a13dc0cf9d977207/grh-compact/src/New/Rows/StackedDeterminants.lean\#L244-L245}{{\leanfont genIdele\_apply\_of\_notMem}}}{StackedDeterminants.lean}{244--245}
\crow{reductionZero}{\href{https://github.com/lsernau/grh-lean/blob/72fead102f663d5c7a407964a13dc0cf9d977207/grh-compact/src/New/Rows/StackedDeterminants.lean\#L247-L250}{{\leanfont reduction\_levelModulus\_eq\_zero\_iff}}}{StackedDeterminants.lean}{247--250}
\crow{genResidue}{\href{https://github.com/lsernau/grh-lean/blob/72fead102f663d5c7a407964a13dc0cf9d977207/grh-compact/src/New/Rows/StackedDeterminants.lean\#L256-L256}{{\leanfont DividedSplitting.\allowbreak{}genResidue}}}{StackedDeterminants.lean}{256--256}
\crow{genResidueSmul}{\href{https://github.com/lsernau/grh-lean/blob/72fead102f663d5c7a407964a13dc0cf9d977207/grh-compact/src/New/Rows/StackedDeterminants.lean\#L258-L266}{{\leanfont DividedSplitting.\allowbreak{}genResidue\_smul\_eq\_iff}}}{StackedDeterminants.lean}{258--266}
\crow{stackedDet}{\href{https://github.com/lsernau/grh-lean/blob/72fead102f663d5c7a407964a13dc0cf9d977207/grh-compact/src/New/Rows/StackedDeterminants.lean\#L277-L277}{{\leanfont stackedDet}}}{StackedDeterminants.lean}{277--277}
\crow{unitModulus}{\href{https://github.com/lsernau/grh-lean/blob/72fead102f663d5c7a407964a13dc0cf9d977207/grh-compact/src/New/Rows/StackedDeterminants.lean\#L280-L280}{{\leanfont unitModulus}}}{StackedDeterminants.lean}{280--280}
\crow{carrierForm}{\href{https://github.com/lsernau/grh-lean/blob/72fead102f663d5c7a407964a13dc0cf9d977207/grh-compact/src/New/Rows/StackedDeterminants.lean\#L289-L291}{{\leanfont carrierForm}}}{StackedDeterminants.lean}{289--291}
\crow{carrierXi}{\href{https://github.com/lsernau/grh-lean/blob/72fead102f663d5c7a407964a13dc0cf9d977207/grh-compact/src/New/Rows/StackedDeterminants.lean\#L297-L298}{{\leanfont carrierXi}}}{StackedDeterminants.lean}{297--298}
\crow{detFromCols}{\href{https://github.com/lsernau/grh-lean/blob/72fead102f663d5c7a407964a13dc0cf9d977207/grh-compact/src/New/Rows/StackedDeterminants.lean\#L304-L313}{{\leanfont exists\_det\_fromCols}}}{StackedDeterminants.lean}{304--313}
\crow{twoSidedSubst}{\href{https://github.com/lsernau/grh-lean/blob/72fead102f663d5c7a407964a13dc0cf9d977207/grh-compact/src/New/Rows/StackedDeterminants.lean\#L317-L318}{{\leanfont twoSidedSubst}}}{StackedDeterminants.lean}{317--318}
\crow{inFirstCols}{\href{https://github.com/lsernau/grh-lean/blob/72fead102f663d5c7a407964a13dc0cf9d977207/grh-compact/src/New/Rows/StackedDeterminants.lean\#L322-L323}{{\leanfont inFirstCols}}}{StackedDeterminants.lean}{322--323}
\crow{stackedDetProd}{\href{https://github.com/lsernau/grh-lean/blob/72fead102f663d5c7a407964a13dc0cf9d977207/grh-compact/src/New/Rows/StackedDeterminants.lean\#L327-L327}{{\leanfont stackedDetProd}}}{StackedDeterminants.lean}{327--327}
\crow{substStackedDetProd}{\href{https://github.com/lsernau/grh-lean/blob/72fead102f663d5c7a407964a13dc0cf9d977207/grh-compact/src/New/Rows/StackedDeterminants.lean\#L332-L337}{{\leanfont twoSidedSubst\_stackedDetProd}}}{StackedDeterminants.lean}{332--337}
\crow{firstColsCoord}{\href{https://github.com/lsernau/grh-lean/blob/72fead102f663d5c7a407964a13dc0cf9d977207/grh-compact/src/New/Rows/StackedDeterminants.lean\#L342-L343}{{\leanfont firstColsCoord}}}{StackedDeterminants.lean}{342--343}
\crow{substFirstColsCoord}{\href{https://github.com/lsernau/grh-lean/blob/72fead102f663d5c7a407964a13dc0cf9d977207/grh-compact/src/New/Rows/StackedDeterminants.lean\#L345-L353}{{\leanfont twoSidedSubst\_firstColsCoord}}}{StackedDeterminants.lean}{345--353}
\crow{blockSigns}{\href{https://github.com/lsernau/grh-lean/blob/72fead102f663d5c7a407964a13dc0cf9d977207/grh-compact/src/New/Rows/RealMatching.lean\#L20-L20}{{\leanfont blockSigns}}}{RealMatching.lean}{20--20}
\crow{posBlockEmb}{\href{https://github.com/lsernau/grh-lean/blob/72fead102f663d5c7a407964a13dc0cf9d977207/grh-compact/src/New/Rows/RealMatching.lean\#L24-L25}{{\leanfont posBlockEmb}}}{RealMatching.lean}{24--25}
\crow{posBlockEmbSpec}{\href{https://github.com/lsernau/grh-lean/blob/72fead102f663d5c7a407964a13dc0cf9d977207/grh-compact/src/New/Rows/RealMatching.lean\#L28-L31}{{\leanfont posBlockEmb\_spec}}}{RealMatching.lean}{28--31}
\crow{cardNegBlockSigns}{\href{https://github.com/lsernau/grh-lean/blob/72fead102f663d5c7a407964a13dc0cf9d977207/grh-compact/src/New/Rows/RealMatching.lean\#L33-L34}{{\leanfont card\_neg\_blockSigns}}}{RealMatching.lean}{33--34}
\crow{matchedPoly}{\href{https://github.com/lsernau/grh-lean/blob/72fead102f663d5c7a407964a13dc0cf9d977207/grh-compact/src/New/Rows/RealMatching.lean\#L38-L39}{{\leanfont matchedPoly}}}{RealMatching.lean}{38--39}
\crow{posFockVector}{\href{https://github.com/lsernau/grh-lean/blob/72fead102f663d5c7a407964a13dc0cf9d977207/grh-compact/src/New/Rows/RealMatching.lean\#L43-L44}{{\leanfont posFockVector}}}{RealMatching.lean}{43--44}
\crow{realWeilStackedDetProd}{\href{https://github.com/lsernau/grh-lean/blob/72fead102f663d5c7a407964a13dc0cf9d977207/grh-compact/src/New/Rows/RealMatching.lean\#L52-L56}{{\leanfont realPairWeil\_stackedDetProd}}}{RealMatching.lean}{52--56}
\crow{realmatching}{\href{https://github.com/lsernau/grh-lean/blob/72fead102f663d5c7a407964a13dc0cf9d977207/grh-compact/src/New/Rows/RealMatching.lean\#L61-L113}{{\leanfont real\_matching}}}{RealMatching.lean}{61--113}
\crow{isRapidModKernel}{\href{https://github.com/lsernau/grh-lean/blob/72fead102f663d5c7a407964a13dc0cf9d977207/grh-compact/src/New/Rows/Carrier.lean\#L42-L45}{{\leanfont isRapid\_modKernel}}}{Carrier.lean}{42--45}
\crow{explicitMatchingOfSum}{\href{https://github.com/lsernau/grh-lean/blob/72fead102f663d5c7a407964a13dc0cf9d977207/grh-compact/src/New/Rows/Carrier.lean\#L52-L75}{{\leanfont explicit\_matching\_of\_sum}}}{Carrier.lean}{52--75}
\crow{classSumRowStabilizer}{\href{https://github.com/lsernau/grh-lean/blob/72fead102f663d5c7a407964a13dc0cf9d977207/grh-compact/src/New/Rows/Carrier.lean\#L80-L103}{{\leanfont stdClassSum\_rowStabilizer\_eq}}}{Carrier.lean}{80--103}
\crow{explicitmatching}{\href{https://github.com/lsernau/grh-lean/blob/72fead102f663d5c7a407964a13dc0cf9d977207/grh-compact/src/New/Rows/Carrier.lean\#L112-L146}{{\leanfont explicit\_matching}}}{Carrier.lean}{112--146}
\crow{classSumDet}{\href{https://github.com/lsernau/grh-lean/blob/72fead102f663d5c7a407964a13dc0cf9d977207/grh-compact/src/New/Rows/Carrier.lean\#L159-L171}{{\leanfont stdClassSum\_det\_modKernel\_eq}}}{Carrier.lean}{159--171}
\crow{carrierIdentity}{\href{https://github.com/lsernau/grh-lean/blob/72fead102f663d5c7a407964a13dc0cf9d977207/grh-compact/src/New/Rows/Carrier.lean\#L178-L199}{{\leanfont stdClassSum\_carrierXi\_eq\_carrierForm}}}{Carrier.lean}{178--199}}
% the notations of the defined functions: one slot for each explicit argument of the Lean declaration
\newcommand{\FixesColsOff}[2]{\operatorname{Fix}_{#1}(#2)}
\newcommand{\rowStabilizer}[3]{\operatorname{Stab}(#1,#2,#3)}
\newcommand{\columnDatum}[5]{\operatorname{cd}_{#1,#2,#3,#4}(#5)}
\newcommand{\HasWeightNormLe}[4]{\operatorname{WN}_{#1}(#2,#3,#4)}
\newcommand{\ChosenGL}[2]{\mathsf{G}(#1,#2)}
\newcommand{\firstBlockHom}[3]{\operatorname{blk}_{#1,#2}(#3)}
\newcommand{\auxLevel}[2]{\mathsf{H}(#1,#2)}
\newcommand{\AuxPair}[3]{\operatorname{Aux}(#1;#2,#3)}
\newcommand{\levelGroup}[1]{\Gamma_{#1}}
\newcommand{\resDetDatum}[2]{\operatorname{res}_{#1}(#2)}
\newcommand{\auxKappa}[1]{\kappa_{#1}}
\newcommand{\detExp}[2]{\epsilon_{#1}(#2)}
\newcommand{\auxLevelChar}[2]{\bar\chi_{#1}(#2)}
\newcommand{\levelChar}[1]{\chi_{#1}}
\newcommand{\FormV}[1]{\mathcal{F}(#1)}
\newcommand{\chosenGLRep}[2]{\operatorname{rep}_{#1}(#2)}
\newcommand{\formAction}[1]{\operatorname{act}_{#1}}
\newcommand{\tensorPoly}[1]{\operatorname{tp}_{#1}}
\newcommand{\harmonicPoly}[1]{P^{#1}}
\newcommand{\auxA}[1]{\mathbf{a}_{#1}}
\newcommand{\auxB}[1]{\mathbf{b}_{#1}}
\newcommand{\pairThetaKernel}[1]{\mathcal{K}_{#1}}
\newcommand{\auxDatum}[1]{\Upsilon_{#1}}
\newcommand{\fockDatum}[2]{\Upsilon_{#1}[#2]}
\newcommand{\conjKernel}[3]{\bar\theta_{#1}(#2,#3)}
\newcommand{\descent}[3]{\operatorname{lift}_{#1}(#2)(#3)}
\newcommand{\rankDiff}[1]{\operatorname{cork}(#1)}
\newcommand{\testValue}[2]{z_{#1}(#2)}
\newcommand{\chosenModulusSq}[1]{\mathfrak{M}_{#1}^{2}}
\newcommand{\reduceToRzero}[1]{\operatorname{red}_{0}(#1)}
\newcommand{\DividedSplitting}[5]{\operatorname{Div}(#1,#2,#3;#4,#5)}
\newcommand{\genIdele}[1]{\varpi_{#1}}
\newcommand{\genIntAdele}[1]{\hat\varpi_{#1}}
\newcommand{\genResidue}[1]{\bar\varpi_{#1}}
\newcommand{\stackedDet}[3]{\Delta_{#1}(#2;#3)}
\newcommand{\unitModulus}[1]{(1)_{#1}}
\newcommand{\carrierXi}[3]{\Xi^{\circ}_{#1,#2}(#3)}
\newcommand{\firstColsCoord}[3]{\operatorname{coord}_{#1}(#2)_{#3}}
\newcommand{\blockSigns}[2]{\operatorname{sgn}_{#1}(#2)}
\newcommand{\posBlockEmb}[2]{\operatorname{pos}_{#1}(#2)}
\newcommand{\posFockVector}[1]{\operatorname{fv}_{#1}}
\newcommand{\UnitMembers}{219}
\newcommand{\UnitInNote}{73}
\newcommand{\UnitProofs}{85}
\newcommand{\UnitSteps}{634}
\newcommand{\UnitAllSteps}{680}


\linespread{1.08}                       % 11 pt on 14.7 pt
\setlength{\parindent}{0pt}
\setlength{\parskip}{5.5pt plus 1.5pt minus 0.5pt}
\setlength{\abovedisplayskip}{7.5pt plus 2pt minus 3pt}
\setlength{\belowdisplayskip}{7.5pt plus 2pt minus 3pt}
\setlength{\abovedisplayshortskip}{3pt plus 2pt}
\setlength{\belowdisplayshortskip}{6pt plus 2pt minus 2pt}
\emergencystretch=1.5em
\clubpenalty=10000 \widowpenalty=10000 \displaywidowpenalty=10000
\relpenalty=7000 \binoppenalty=9000
\raggedbottom
\pagestyle{fancy}\fancyhf{}\renewcommand{\headrulewidth}{0pt}
\fancyfoot[C]{\sffamily\fontsize{7.6}{9}\selectfont\color{inkfaint}\thepage\,/\,\pageref*{LastPage}}
\setlist{topsep=2pt,partopsep=0pt,parsep=0pt,itemsep=4pt,leftmargin=2.1em,labelsep=0.55em}

\newcommand{\Q}{\mathbb{Q}}
\newcommand{\Z}{\mathbb{Z}}
\newcommand{\F}{\mathbb{F}}
\newcommand{\len}{\operatorname{len}}
\newcommand{\cok}{\operatorname{cok}}
\newcommand{\End}{\operatorname{End}}
\newcommand{\Hom}{\operatorname{Hom}}
\newcommand{\Gal}{\operatorname{Gal}}
\newcommand{\diag}{\operatorname{diag}}
\newcommand{\kw}[1]{\emph{#1}}
\newcommand{\m}[1]{\mbox{$#1$}}        % a short formula that must not be broken across lines
\newcommand{\head}[1]{\par\addvspace{17pt plus 4pt minus 2pt}\Needspace{6\baselineskip}%
  \noindent{\fontsize{13.2}{16}\selectfont\bfseries\color{ink}#1}\par\nopagebreak\vspace{-1pt}}
\newcommand{\sect}[2]{\refstepcounter{section}\label{#2}\par\addvspace{19pt plus 4pt minus 2pt}\Needspace{7\baselineskip}%
  \noindent\llap{{\fontsize{13.2}{16}\selectfont\bfseries\color{inkfaint}\thesection}\hspace{9pt}}%
  {\fontsize{13.2}{16}\selectfont\bfseries\color{ink}#1}\par\nopagebreak\vspace{-1pt}}
\newcommand{\subhead}[1]{\par\addvspace{9pt plus 2pt}\Needspace{4\baselineskip}\noindent\textbf{#1}\enspace\ignorespaces}
% where a statement lives in the Lean archive
\newcommand{\leantag}{{\sffamily\bfseries\fontsize{6.3}{8}\selectfont\color{leancol}\addfontfeatures{LetterSpace=7}LEAN}}
\newcommand{\sep}{\unskip\nobreak\hspace{5pt}{\color{inkfaint}\textperiodcentered}\hspace{5pt plus 2pt}\ignorespaces}
\newcommand{\leanline}[1]{\par\nopagebreak\vspace{-2.5pt}{\fontsize{8.9}{12.4}\selectfont\color{inksoft}\raggedright
  \noindent\llap{\leantag\hspace{9pt}}#1\par}\vspace{1pt}}
% a displayed statement in a tinted box: #1 colour, #2 text
\newcommand{\statement}[2]{\par\addvspace{7pt}\begingroup\setlength{\fboxsep}{8.5pt}\setlength{\fboxrule}{0.6pt}%
  \noindent\fcolorbox{#1!30!white}{#1!3!white}{\parbox{\dimexpr\linewidth-2\fboxsep-2\fboxrule\relax}{#2}}%
  \endgroup\par\addvspace{7pt}}
\newcommand{\nt}[2]{\par\noindent\hangindent=1.1em\hangafter=1 #1\enspace{\color{inksoft}#2}\par\vspace{1pt}}
\newcommand{\legdot}[1]{\tikz[baseline=-0.62ex]\fill[#1] (0,0) circle (2.5pt);\hspace{3pt}}
\newcommand{\crow}[4]{\csname CS@#1\endcsname & #2 & \textcolor{inksoft}{#3} & \textcolor{inksoft}{#4}\\}
\newcommand{\leanquoteline}[2]{\par\noindent\hangindent=\dimexpr#1\fontdimen2\font+2.2em\relax\hangafter=1
  \hspace*{#1\fontdimen2\font}#2}
\newcommand{\quoted}[1]{\par\vspace{3pt}{\ttfamily\fontsize{7.5}{10.4}\selectfont\color{leancol}\raggedright
  \setlength{\parskip}{0pt}\leanquote{#1}\par}\vspace{4pt}}
\newcommand{\bx}[3]{{\sffamily\bfseries\fontsize{7.3}{9.2}\selectfont\color{#1!82!black}#2}\\[1.5pt]#3}
\newcommand{\cs}[2]{\expandafter\newcommand\csname CS@#1\endcsname{#2}}

% theorems
\newtheoremstyle{plainx}{7pt plus 2pt}{7pt plus 2pt}{\itshape}{}{\bfseries}{.}{0.5em}{}
\newtheoremstyle{uprightx}{7pt plus 2pt}{7pt plus 2pt}{}{}{\bfseries}{.}{0.5em}{}
\theoremstyle{plainx}
\newtheorem{theorem}{Theorem}[section]
\newtheorem{lemma}[theorem]{Lemma}
\newtheorem{proposition}[theorem]{Proposition}
\newtheorem{corollary}[theorem]{Corollary}
\theoremstyle{uprightx}
\newtheorem{definition}[theorem]{Definition}
\newtheorem{setting}[theorem]{Setting}          % hypotheses shared by several statements; each of them cites it by number
\newtheorem{remark}[theorem]{Remark}
\newtheorem{example}[theorem]{Example}
\renewcommand{\qedsymbol}{{\setlength{\fboxsep}{0pt}\setlength{\fboxrule}{0.5pt}\fbox{\rule{0pt}{0.46em}\hspace{0.46em}}}}
% what holds throughout the note: the ambient objects only (implicit and instance arguments), never a hypothesis
\newenvironment{conventions}{\par\addvspace{5pt}\noindent\textbf{Conventions.}\enspace\ignorespaces}{\par\addvspace{3pt}}
\makeatletter
\renewenvironment{proof}[1][Proof]{\par\pushQED{\qed}\normalfont\topsep0pt\partopsep0pt
  \trivlist\item[\hskip\labelsep\itshape #1.]\ignorespaces}{\popQED\endtrivlist\@endpefalse}
\makeatother

% where each declaration of the concordance is treated: \cs{key}{place}, the place naming a \label or a lettered item
\cs{FixesColsOff}{Definition~\ref{d:cols}}
\cs{rowStabilizer}{Definition~\ref{d:cols}}
\cs{columnDatum}{Definition~\ref{d:cols}}
\cs{columnDatumMul}{Lemma~\ref{l:cols}}
\cs{HasWeightNormLe}{Definition~\ref{d:wn}}
\cs{HasWeightNormLeProd}{Lemma~\ref{l:wnprod}}
\cs{ChosenGL}{Definition~\ref{d:levels}}
\cs{firstBlockHom}{Definition~\ref{d:levels}}
\cs{auxLevel}{Definition~\ref{d:levels}}
\cs{AuxPair}{Definition~\ref{d:aux}}
\cs{levelGroup}{Definition~\ref{d:auxobj}}
\cs{resDetDatum}{Definition~\ref{d:auxobj}}
\cs{auxKappa}{Definition~\ref{d:auxobj}}
\cs{detExp}{Definition~\ref{d:auxobj}}
\cs{auxLevelChar}{Definition~\ref{d:auxobj}}
\cs{levelChar}{Definition~\ref{d:auxobj}}
\cs{FormV}{Definition~\ref{d:auxobj}}
\cs{chosenGLRep}{Definition~\ref{d:auxobj}}
\cs{formAction}{Definition~\ref{d:auxobj}}
\cs{tensorPoly}{Definition~\ref{d:auxobj}}
\cs{tensorPolyVec}{Remark~\ref{r:auxfacts}}
\cs{harmonicPoly}{Definition~\ref{d:auxobj}}
\cs{auxA}{Definition~\ref{d:auxobj}}
\cs{auxB}{Definition~\ref{d:auxobj}}
\cs{pairThetaKernel}{Definition~\ref{d:auxobj}}
\cs{auxDatum}{Definition~\ref{d:auxobj}}
\cs{fockDatum}{Definition~\ref{d:auxobj}}
\cs{conjKernel}{Definition~\ref{d:auxobj}}
\cs{descent}{Definition~\ref{d:auxobj}}
\cs{heckePolyOnForms}{Definition~\ref{d:auxobj}}
\cs{heckePolySmul}{Remark~\ref{r:auxfacts}}
\cs{rankDiff}{Definition~\ref{d:auxobj}}
\cs{testValue}{Definition~\ref{d:auxobj}}
\cs{Lifts}{Definition~\ref{d:auxobj}}
\cs{chosenModulusSq}{Definition~\ref{d:moduli}}
\cs{mapBlock}{Remark~\ref{r:book}}
\cs{levelModulusIdealLe}{Remark~\ref{r:book}}
\cs{placesLevelModulus}{Remark~\ref{r:book}}
\cs{reduceToRzero}{Definition~\ref{d:moduli}}
\cs{liftModSupported}{Remark~\ref{r:book}}
\cs{atPlaceTwist}{Remark~\ref{r:book}}
\cs{DividedSplitting}{Definition~\ref{d:div}}
\cs{genIdele}{Definition~\ref{d:div}}
\cs{genIntAdele}{Definition~\ref{d:div}}
\cs{genIdeleApply}{Remark~\ref{r:book}}
\cs{reductionZero}{Lemma~\ref{l:red}}
\cs{genResidue}{Definition~\ref{d:div}}
\cs{genResidueSmul}{Lemma~\ref{l:genres}}
\cs{stackedDet}{Definition~\ref{d:stacked}}
\cs{unitModulus}{Definition~\ref{d:moduli}}
\cs{carrierForm}{Definition~\ref{d:carrier}}
\cs{carrierXi}{Definition~\ref{d:carrier}}
\cs{detFromCols}{Lemma~\ref{l:detcols}}
\cs{twoSidedSubst}{Definition~\ref{d:stacked}}
\cs{inFirstCols}{Definition~\ref{d:stacked}}
\cs{stackedDetProd}{Definition~\ref{d:stacked}}
\cs{substStackedDetProd}{Lemma~\ref{l:subst}}
\cs{firstColsCoord}{Definition~\ref{d:stacked}}
\cs{substFirstColsCoord}{Lemma~\ref{l:coord}}
\cs{blockSigns}{Definition~\ref{d:real}}
\cs{posBlockEmb}{Definition~\ref{d:real}}
\cs{posBlockEmbSpec}{Remark~\ref{r:pos}}
\cs{cardNegBlockSigns}{Remark~\ref{r:pos}}
\cs{matchedPoly}{Definition~\ref{d:real}}
\cs{posFockVector}{Definition~\ref{d:real}}
\cs{realWeilStackedDetProd}{Lemma~\ref{l:realweil}}
\cs{realmatching}{Proposition~\ref{p:real}}
\cs{isRapidModKernel}{Lemma~\ref{l:rapid}}
\cs{explicitMatchingOfSum}{Proposition~\ref{p:sum}}
\cs{classSumRowStabilizer}{Lemma~\ref{l:stab}}
\cs{explicitmatching}{Theorem A}
\cs{classSumDet}{Theorem B}
\cs{carrierIdentity}{Theorem C}
\newcommand{\CC}{\mathbb{C}}
\newcommand{\ZZ}{\mathbb{Z}}
\newcommand{\NN}{\mathbb{N}}
\newcommand{\AAA}{\mathbb{A}}
\newcommand{\ix}[1]{\mathsf{#1}}
\newcommand{\OEr}{\mathcal{O}_E}
\newcommand{\Ohat}{\widehat{\mathcal{O}}_E}
\newcommand{\UU}{\mathrm{U}}
\newcommand{\GL}{\mathrm{GL}}
\newcommand{\Mat}{\mathrm{M}}
\newcommand{\KK}[1]{\mathbf{K}_{#1}}
\newcommand{\ip}[2]{\langle #1,#2\rangle}
\newcommand{\clsum}[2]{\sideset{}{^{\mathrm{cl}}}\sum_{#1\bmod #2}}
\newcommand{\tk}[7]{\theta^{#1,#2}_{#3,#4;#5}[#6,#7]}
\newcommand{\LP}[3]{\Pi_{#1}(#2,#3)}
\newcommand{\carrier}[2]{\mathfrak{C}_{#1,#2}}
\newcommand{\sub}[2]{\operatorname{sub}_{#1,#2}}
\newcommand{\sdp}[2]{\Delta(#1;#2)}
\newcommand{\mpoly}[3]{Q_{#1,#2}(#3)}
\newcommand{\infc}{\operatorname{col}}
\newcommand{\inputbox}[1]{\statement{inksoft}{#1}}
\newcommand{\transp}{^{\mathsf{T}}}
\newcommand{\herm}{^{\mathsf{H}}}
% letters that the coordinator may change in one place
\newcommand{\dL}{\check{\Psi}}         % the partner on the large side (Lean Φ; the letter Φ is taken by Φ_{d,M})
\newcommand{\Plc}{\mathrm{Pl}}         % a finite set, in the applications the infinite places (Lean Φ)
\newcommand{\lpl}{\Sigma}              % level places (Lean S)

\begin{document}
\color{ink}

{\raggedright
{\sffamily\bfseries\fontsize{7.2}{9}\selectfont\color{leancol}\addfontfeatures{LetterSpace=8}READBACK NOTE · CHAPTER VII · UNIT 143 · PART 2 OF 4\par}
\vspace{5pt}
{\fontsize{21.5}{25.5}\selectfont\bfseries Auxiliary dual pairs, determinant polynomials and the carrier: the explicit
matching of theta kernels\par}
\vspace{5pt}
{\fontsize{12.4}{16}\selectfont\itshape\color{inksoft} One unit of the Lean artifact \ArtifactName, read back into
mathematics\par}}
\vspace{8pt}
{\color{rulegrey}\hrule height 0.6pt}
\vspace{7pt}
\thispagestyle{empty}

\textbf{Summary.} This note sets up the objects with which the rows of the proof are compared with definite
groups: the record of an \emph{auxiliary dual pair}, with its level, its space of algebraic forms, its theta descent
and its Hecke test; the determinants of a constant matrix stacked over a matrix of variables; and the \emph{carrier},
a fixed holomorphic theta lift of the constant~$1$. It then proves three identities between theta kernels. The
\emph{explicit matching} (Theorem~A) says that a theta kernel summed over the classes of a definite unitary group
equals a theta kernel integrated over a unitary group of smaller rank, with explicit determinant polynomials at the
infinite places on both sides, up to a constant whose absolute value is given by Hecke $L$-values. Two consequences
follow: a class sum of a kernel of two definite groups of equal rank is a constant of this kind (Theorem~B), and a
twisted class sum against an explicit form is a constant multiple of the carrier (Theorem~C), so that a twist
survives only as a number.

\statement{leancol}{\fontsize{9.6}{13.2}\selectfont\textbf{What this note is.} The artifact \ArtifactName{} is a
body of Lean code: the compact edition of a formal proof, due to L.~Alpöge with the AI system Claude, of the Riemann
hypothesis and of its generalisation to Dirichlet $L$-functions. That headline is a claim. The authors report that
every file compiles with Lean v4.35.0 against a pinned version of Mathlib and that Lean's kernel accepted the
proofs; an independent build is recorded in the handout's record of checks when it reports. No mathematician has
read the argument through, and nothing here says that the Riemann hypothesis is proved. This note is the second of
four on one unit of the artifact, the folder of the rows; it treats \UnitInNote{} of the unit's \UnitMembers{}
declarations, those of five files. All definitions are given. The lemmas of linear algebra,
Proposition~\ref{p:real} and Theorems A and~B are proved in full from results of the background library and from
one result of another unit, the exchange theorem (Inputs 1 to~9); these are stated and not proved here. The proofs
of Proposition~\ref{p:sum} and of Theorem~C, which rests on it, each describe one computation without carrying it
out: the exact form of the translation rule in the first, the bookkeeping of the characters in the second. They say
so where they stand. The archimedean step, Proposition~\ref{p:real}, was read line by
line in the Lean text. Every statement below was compared with its Lean text, argument by argument. The Lean text
read is \PinText; Appendix~A is the concordance of \ConcDecls{} declarations.}

\sect{The statements}{s:intro}
The objects are defined in §\ref{s:setting}: the theta kernel \m{\tk{S}{T}{\alpha}{\beta}{\mathfrak{m}}{\Psi}{P}(g,h)}
of two hermitian spaces, the class sum \m{\clsum{h}{\KK{\ix a}}} over a definite unitary group, tame characters and
the product \m{\LP{\mathfrak{m}}{\nu}{d}} of $L$-values in Setting~\ref{set:bg}; the determinant polynomials
\m{\sdp{e}{\eta}} and \m{\mpoly{f}{j}{\eta}} in Definitions \ref{d:stacked} and~\ref{d:real}; the carrier in
Definition~\ref{d:carrier}.

\statement{arith}{\phantomsection\label{t:A}\textbf{Theorem A} (explicit matching).\enspace
\emph{In Setting~\ref{set:bg}, let \m{\ix n,\ix m,\ix a,\ix b,\ix p,\ix l} be finite index sets, $s$ a sign vector on~$\ix n$
and $W$ a hermitian space over~$E$ with index set~$\ix n$ and signs~$s$. Let \m{f\colon\ix p\sqcup\ix l\to\ix b} be a
bijection and $V'$ a hermitian space with index set~$\ix b$ and signs \m{\blockSigns{f}{\cdot}}. Let
\m{\sigma\colon\ix b\sqcup\ix a\to\ix n}, \m{e\colon\ix p\sqcup\ix a\to\ix m} and \m{j\colon\ix m\sqcup\ix l\to\ix n}
be bijections with \m{s=\blockSigns{j}{\cdot}}, let \m{d\ge0} be an integer with \m{|\ix a|=|\ix b|+2d}, suppose
\m{2|\ix l|+2\le|\ix n|}, and let \m{\iota=(\iota_\tau)} be a family of finite sets indexed by the infinite places
of~$E$. Then there is a real number~$C$ such that for every modulus~$\mathfrak{m}$, every character~$\alpha$, weight~$\gamma$
and character~$\chi$ of \m{\OEr/\mathfrak{m}} for which $\alpha$ is tame for \m{(\mathfrak{m},\gamma,\chi)} and
\m{\gamma_\tau=|\ix l|-|\iota_\tau|} for all~$\tau$, there is \m{c\in\CC} with
\[|c|^2=\frac{C}{\LP{\mathfrak{m}}{\alpha^{-1}\upsilon^{|\ix b|+d}}{d}}\]
such that for every family \m{\eta=(\eta_{\tau k})} of complex \m{\ix p\times\ix m} matrices, \m{k\in\iota_\tau}, for all
partners $\dL$, $\Psi$ for \m{(\sigma,\chi^{-1})}, every character~$\beta$ and every \m{g\in\UU(W)(\AAA)}:
\[c\cdot\clsum{h}{\KK{\ix a}}\tk{W}{I_{\ix a}}{\alpha}{\beta}{\mathfrak{m}}{\dL}{\sdp{e}{\eta}^{j}}(g,h)
=\int_{[\UU(V')]}\tk{V'}{W}{\alpha/\beta}{\alpha}{\mathfrak{m}}{\Psi}{\mpoly{f}{j}{\eta}}(x,g)\,dx .\]
Here \m{\sdp{e}{\eta}^{j}} is the polynomial \m{\sdp{e}{\eta}} with the column index \m{c\in\ix m} of its variables
replaced by \m{j(c)\in\ix n}.}}
\leanline{\lean{explicitmatching}}

\statement{arith}{\phantomsection\label{t:B}\textbf{Theorem B} (a class sum of a theta kernel as Hecke values).\enspace
\emph{In Setting~\ref{set:bg}, let \m{m\ge1} be an even integer and let $c$ assign an integer \m{c_\tau\ge0} to each
infinite place~$\tau$ of~$E$. Then there is a real number~$C$ such that for every modulus~$\mathfrak{m}$, every
character~$\alpha$ and every character~$\chi$ of \m{\OEr/\mathfrak{m}} with $\alpha$ tame for
\m{(\mathfrak{m},-c,\chi)} there is \m{\gamma\in\CC} with
\m{|\gamma|^2=C/\LP{\mathfrak{m}}{\alpha^{-1}\upsilon^{m/2}}{m/2}} and, for every \m{g\in\UU(I_m)(\AAA)},
\[\gamma\cdot\clsum{h}{\KK{m}}\tk{I_m}{I_m}{\alpha}{\alpha}{\mathfrak{m}}{\chi\circ\det}{\textstyle\prod_\tau\det(x_\tau)^{c_\tau}}(g,h)=1 .\]}

\vspace{4pt}\phantomsection\label{t:C}\textbf{Theorem C} (the carrier: one fixed form under all twists).\enspace
\emph{In Setting~\ref{set:bg}, let \m{m\ge1} be an even integer. Then there is a real number~$C$ such that for every
modulus~$\mathfrak{m}$, every character~$c$ and every character~$\chi$ of \m{\OEr/\mathfrak{m}} with $c$ tame for
\m{(\mathfrak{m},0,\chi)} there is \m{\gamma\in\CC} with
\[|\gamma|^2=\frac{C}{\LP{\mathfrak{m}}{c^{-1}\upsilon^{m}}{m}\cdot\LP{\mathfrak{m}}{c\,\upsilon^{m/2}}{m/2}}\]
and, for every \m{x\in\UU(I_{2m,m})(\AAA)},
\[\gamma\cdot\clsum{h}{\KK{3m}}\carrierXi{m}{c}{h}\cdot\tk{I_{2m,m}}{I_{3m}}{\upsilon^{m}}{c}{\mathfrak{m}}{\chi\circ\det}{1}(x,h)=\carrier{E}{m}(x) .\]}}
\leanline{\lean{classSumDet}\sep \lean{carrierIdentity}\sep in Theorems B and~C the integer~$m$ is the Lean
variable; it is not the number of lifts of the data, which is its value where the theorems are applied}

In Theorem~C the left side depends on the character~$c$, which in the applications carries the twist, and the right
side does not: the twist is seen only through the number~$\gamma$, of known absolute value. In Theorems A to~C the
constant~$C$ may be zero; then the corresponding number is zero and the identity says that the right side
vanishes. That the carrier is not zero is proved in the note \emph{The height of a ratio of lift volumes}.

\sect{The setting}{s:setting}
\begin{conventions}\phantomsection\label{c:conventions}
Finite index sets are written in sans-serif, with the letters of the formal text: \m{\ix r,\ix n,\ix N,\ix m,\ix a,\ix b,\ix c,
\ix p,\ix l,\ix L,\ix k}; \m{|\ix n|} is the number of elements of~$\ix n$, and an \m{\ix n\times\ix r} matrix has rows
indexed by~$\ix n$ and columns by~$\ix r$. For a matrix~$z$, \m{z\transp} is its transpose and \m{z\herm} its conjugate
transpose. The ambient objects of the declarations of this note, their implicit and instance arguments, are: the
commutative rings $R$ and~$R'$ of §\ref{s:cols}; the finite set~$\Plc$ of places and the finite index sets of each
statement; and, from Definition~\ref{d:aux} on, a field~$E$ that is a CM field with a CM type and with Weil data, as
in Setting~\ref{set:bg}. Each statement names the objects it uses. In the proofs, when a statement of the formal text
has implicit arguments that a later hypothesis determines (a modulus~$\mathfrak{m}$, characters, signs, a
polynomial), they are introduced with that hypothesis.
\end{conventions}

\begin{setting}[the background: fields, characters, spaces, kernels, weights]\label{set:bg}
All of this is defined in the background library of the artifact; the catalogue of the handout gives references.

\emph{(a) The field.} $E$ is a CM field with a CM type (for each infinite place~$\tau$ a complex embedding that
induces it) and with \kw{Weil data}: the choices that fix the Weil representation, among them a character~$\upsilon$ of
the norm-one ideles of~$E$, the restriction of a splitting character. Finite places are written $w$, with residue
field~$k_w$, norm~$q_w$ and complex conjugate~$\bar w$; $\Ohat$ is the ring of integral finite adeles, $\AAA$ the
adeles and $\AAA_f$ the finite adeles of the totally real subfield.

\emph{(b) Moduli.} A \kw{modulus} is a non-zero ideal~$\mathfrak{m}$ of~$\OEr$ divisible by at most one of $w$,
$\bar w$ for each~$w$. It has a chosen reduction \m{\Ohat\to\OEr/\mathfrak{m}} extending the one of~$\OEr$ and a chosen
\kw{generator}, an integral idele equal to~$1$ at the places not dividing~$\mathfrak{m}$ that generates the kernel of
the reduction. \kw{Level places} are a pair~$\lpl$ of finite families of finite places, the \kw{chosen places} and the
\kw{twist places}, such that all these places are distinct and none is the conjugate of another;
\m{\mathfrak{M}_\lpl} and \m{\mathfrak{Q}_\lpl} are the products of the primes of the two families, so that
\m{\mathfrak{M}_\lpl^{c}\mathfrak{Q}_\lpl} and \m{\mathfrak{M}_\lpl^{c}} are moduli for every integer \m{c\ge0}.

\emph{(c) Characters.} A \kw{character} is a continuous character of the norm-one ideles of~$E$ trivial on the
norm-one elements of~$E$. It has \kw{weight}~$e$, an integer~$e_\tau$ for each~$\tau$, and \kw{tame part}
\m{\psi=(\psi_w)}, a character of \m{k_w^\times} for each~$w$, almost all trivial, if its value at~$u$ is
\m{\prod_\tau u_\tau^{e_\tau}\prod_w\psi_w(u\bmod w)} whenever the finite part of~$u$ is integral. It is \kw{tame for}
\m{(\mathfrak{m},e,\chi)}, for a modulus~$\mathfrak{m}$ and a multiplicative character~$\chi$ of \m{\OEr/\mathfrak{m}},
if \m{\mathfrak{m}\ne(1)}, it has weight~$e$ and a tame part that is non-trivial exactly at the primes
dividing~$\mathfrak{m}$, and $\chi$ is the product of the~$\psi_w$. For a set~$T$ of finite places, a
character~$\nu$ and an integer \m{d\ge0}, with \m{L^{T}(z,\nu)} the partial Hecke $L$-function of the base change
\m{x\mapsto\nu(x/\bar x)} of~$\nu$ to the ideles of~$E$, without its Euler factors at the places of~$T$ and at their
complex conjugates,
\[\LP{T}{\nu}{d}=\Bigl(\prod_{w\in T}q_w\Bigr)^{d^2}\cdot\prod_{z=1}^{d}\bigl|L^{T}(z,\nu)\bigr|^2,\]
and \m{\LP{\mathfrak{m}}{\nu}{d}} is this number for the set of primes dividing~$\mathfrak{m}$.

\emph{(d) Spaces and groups.} For a finite index set~$\ix n$ and a sign vector \m{s\colon\ix n\to\{\pm1\}}, a
\kw{hermitian space with index set~$\ix n$ and signs~$s$} is a hermitian matrix over~$E$ together with an adelic
point that brings it to \m{\operatorname{diag}(s)} at the infinite places and to~$1$ at the finite places. $I_{\ix n}$
is the standard positive definite one, $I_m$ the one with index set \m{\{1,\dots,m\}}, \m{I_{2m,m}} the diagonal
space with $2m$ signs~$+1$ and $m$ signs~$-1$, $S^{-}$ the space~$S$ with the opposite form and \m{S\oplus T} the
orthogonal sum. \m{\UU(S)(\AAA)} is the adelic unitary group, \m{[\UU(S)]} its automorphic quotient with its invariant
measure, and \m{G_\infty(s)} the product over the infinite places of the real unitary groups of \m{\operatorname{diag}(s)}.
For the definite space, \m{\KK{\ix n}=\UU(I_{\ix n})(\Ohat)} is the stabiliser of the standard lattice, and a
\kw{level} is a subgroup of finite index in an open subgroup of it. For a level~$\Gamma$ and a function~$F$ on
\m{\UU(I_{\ix n})(\AAA_f)}, the \kw{class sum} is
\[\clsum{h}{\Gamma}F(h)=\sum_{\mathfrak{c}}\frac{F(h_{\mathfrak{c}})}{\#\bigl(\UU(I_{\ix n})(E)\cap h_{\mathfrak{c}}\Gamma h_{\mathfrak{c}}^{-1}\bigr)},\]
the sum over the double cosets~$\mathfrak{c}$ of \m{\UU(I_{\ix n})(E)\backslash\UU(I_{\ix n})(\AAA_f)/\Gamma} with chosen
representatives~$h_{\mathfrak{c}}$.

\emph{(e) Theta kernels.} Let $S$ and~$T$ be hermitian spaces with index sets $\ix n$ and~$\ix n'$ and $\alpha$,
$\beta$ characters. The Weil representation \m{\omega_{\alpha,\beta}} of \m{\UU(S)(\AAA)\times\UU(T)(\AAA)} acts on
the Schwartz functions on \m{\operatorname{Hom}(S,T)}. At the infinite places it is realised in a Fock model: a
complex polynomial~$P$ in the variables \m{x_{\tau,i,c}}, with $\tau$ an infinite place, \m{i\in\ix n'} and
\m{c\in\ix n}, has a \kw{Fock vector}~$v_P$, the Fock vectors are dense in the Fock space, and the groups
\m{G_\infty} act through the \kw{real Weil representation} \m{\omega^{\infty}_{\gamma,\gamma'}} with integer
weights $\gamma$, $\gamma'$. For a modulus~$\mathfrak{m}$, a complex function~$\Psi$ on the \m{\ix n'\times\ix n}
matrices over \m{\OEr/\mathfrak{m}} and a polynomial~$P$, the \kw{theta kernel} is
\[\tk{S}{T}{\alpha}{\beta}{\mathfrak{m}}{\Psi}{P}(g,h)=\theta\bigl(\omega_{\alpha,\beta}(g,h)\,(v_P\otimes\Psi)\bigr),\qquad g\in\UU(S)(\AAA),\ h\in\UU(T)(\AAA),\]
where $\theta$ is the theta functional and \m{v_P\otimes\Psi} the Schwartz function with archimedean part~$v_P$
and finite part the function~$\Psi$ of the reduction modulo~$\mathfrak{m}$ at the primes dividing~$\mathfrak{m}$,
applied to the vacuum vector elsewhere. When \m{T=I_{\ix n'}} and $h$ is a finite adelic point, the same notation
is used with the infinite component of~$h$ equal to~$1$. Functions $\dL$ on \m{\ix a\times\ix n} matrices and $\Psi$
on \m{\ix n\times\ix b} matrices over a commutative ring are \kw{partners for} \m{(\sigma,\chi)}, with
\m{\sigma\colon\ix b\sqcup\ix a\to\ix n} a bijection and $\chi$ a multiplicative character of the ring, if $\dL$
vanishes off the matrices formed by the rows \m{\sigma(\ix a)} of~$k^{-1}$ and $\Psi$ off those formed by the columns
\m{\sigma(\ix b)} of~$k$, for invertible \m{\ix n\times\ix n} matrices~$k$, and the value of~$\dL$ at the first is
\m{\chi(\det k)} times the value of~$\Psi$ at the second.

\emph{(f) Weights and polynomials.} $\Plc$ is a finite set, in the applications the set of infinite places. A
\kw{weight}~$\kappa$ assigns to each \m{\tau\in\Plc} and each height \m{h\ge0} a number \m{\kappa_\tau(h)\ge0} of
columns. For a finite index set~$\ix n$, the \kw{standard fillings} of weight~$\kappa$ with rows~$\ix n$ index a
basis~$b_l$ of the polynomials in matrix variables that are semi-invariant of weight~$\kappa$ under the Borel
subgroup: products of minors. With the integer Gram matrix~$G_\kappa$ of this basis, a complex coordinate
vector~$\xi$ indexed by the standard fillings has the polynomial
\m{\operatorname{pol}_\kappa(\xi)=\sum_l(G_\kappa\xi)_l\,b_l} and the squared norm
\m{\|\xi\|_\kappa^2=\operatorname{Re}(\xi\herm G_\kappa\xi)}. The \kw{Fock pairing} of two polynomials is
\m{\ip{P}{Q}=P(\partial)Q(0)}; the \kw{hook number}~$h_\kappa$ is the pairing with itself of the product of leading
minors of weight~$\kappa$; the \kw{Schur matrix} \m{S_\kappa(z)} of matrices~$z_\tau$ expresses the polynomials
\m{b_l(zy)} in the basis.
\end{setting}

\subsection*{Matrices with prescribed columns}\label{s:cols}
\begin{definition}[stabiliser and column datum]\label{d:cols}
Let $R$ and~$R'$ be commutative rings, \m{\ix r} and~\m{\ix N} finite index sets, \m{\ix c} an index set,
\m{\chi\colon R\to\CC} a multiplicative character (zero on the non-units), \m{\varphi\colon R\to R'} a ring
homomorphism, $w$ a \m{\ix c\times\ix N} matrix over~$R'$ and \m{e\colon\ix r\hookrightarrow\ix N} an embedding. Let
\m{I_e} be the \m{\ix N\times\ix r} matrix whose column~$i$ is the standard basis vector of index~\m{e(i)}.
\begin{enumerate}[label=\textup{(\roman*)},leftmargin=2.3em,itemsep=2pt]
\item For an \m{\ix N\times\ix N} matrix~$z$ over~$R$, \m{\FixesColsOff{e}{z}} is the proposition
\m{z\,(1-I_eI_e\transp)=1-I_eI_e\transp}: the columns of~$z$ with index outside \m{e(\ix r)} are those of the identity
matrix.
\item \m{\rowStabilizer{\varphi}{w}{e}} is the subgroup of the \m{A\in\GL_{\ix N}(R)} with \m{\FixesColsOff{e}{A}} and
\m{w\,\varphi(A)=w}.
\item For an \m{\ix N\times\ix r} matrix~$x$ over~$R$, the \kw{column datum} is
\[\columnDatum{\chi}{\varphi}{w}{e}{x}=\begin{cases}\chi\bigl(\det(I_e\transp x)\bigr)&\text{if }w\,\varphi(x)=w\,I_e,\\ 0&\text{otherwise,}\end{cases}\]
where \m{I_e\transp x} consists of the rows \m{e(\ix r)} of~$x$.
\end{enumerate}
\end{definition}
\leanline{\lean{FixesColsOff}\sep \lean{rowStabilizer}\sep \lean{columnDatum}}

\begin{lemma}[on the stabiliser the datum is a product]\label{l:cols}
Let $R$ and~$R'$ be commutative rings, \m{\ix r} and~\m{\ix N} finite index sets, \m{\ix a} and~\m{\ix c} index sets,
$\chi$ a multiplicative character of~$R$, \m{\varphi\colon R\to R'} a ring homomorphism, $w$ a \m{\ix c\times\ix N}
matrix over~$R'$ and \m{e\colon\ix r\hookrightarrow\ix N} an embedding. Let $F$ be a complex function on the
\m{\ix N\times\ix a} matrices over~$R$ with \m{F(Ay)=F(y)} for all \m{A\in\rowStabilizer{\varphi}{w}{e}} and all~$y$.
Then for every \m{\ix N\times\ix a} matrix~$y$ and every \m{\ix N\times\ix N} matrix~$z$ over~$R$
\[\columnDatum{\chi}{\varphi}{w}{e}{zI_e}\cdot\bigl[\FixesColsOff{e}{z}\bigr]\,F(y)=\bigl[z\in\rowStabilizer{\varphi}{w}{e}\bigr]\,\chi(\det z)\,F(z^{-1}y),\]
where a bracket is $1$ if the proposition in it holds and $0$ if not.
\end{lemma}
\begin{proof}
Suppose \m{\FixesColsOff{e}{z}}. Then \m{z=1+(z-1)I_eI_e\transp}, so
\[\det z=\det\bigl(1+I_e\transp(z-1)I_e\bigr)=\det(I_e\transp zI_e),\]
since \m{I_e\transp I_e=1}; and \m{w\varphi(z)=w} holds if and only if \m{w\varphi(z)I_e=wI_e}, because the columns
outside \m{e(\ix r)} agree on both sides. So \m{\columnDatum{\chi}{\varphi}{w}{e}{zI_e}} is \m{\chi(\det z)} if
\m{w\varphi(z)=w}, and $0$ if not. If $z$ is not invertible, then \m{\chi(\det z)=0} and \m{z} is in no subgroup of
\m{\GL_{\ix N}(R)}: both sides vanish. If $z$ is invertible, the left side is \m{\chi(\det z)F(y)} when $z$ lies in the
stabiliser and $0$ otherwise, and for $z$ in the stabiliser \m{F(y)=F(z^{-1}y)} by the invariance of~$F$
under~$z^{-1}$.
\end{proof}
\leanline{\lean{columnDatumMul}}

\subsection*{Weights and norms of products}
\begin{definition}[weight with a norm bound]\label{d:wn}
Let $\Plc$ and~$\ix n$ be finite, $p$ a proposition, $\kappa$ a weight, $F$ a polynomial and $b$ a real number.
\m{\HasWeightNormLe{p}{\kappa}{F}{b}} means: there is a coordinate vector~$\xi$ of weight~$\kappa$ with rows~$\ix n$
such that \m{\operatorname{pol}_\kappa(\xi)=F} and, if $p$ holds, \m{\|\xi\|_\kappa^2\le b}.
\end{definition}
\leanline{\lean{HasWeightNormLe}}

\begin{lemma}[products]\label{l:wnprod}
Let $\Plc$ and~$\ix n$ be finite, $p$ a proposition, $s$ a finite set and, for \m{i\in s}, let $\kappa_i$ be a
weight, $F_i$ a polynomial and $b_i$ a real number. If \m{\HasWeightNormLe{p}{\kappa_i}{F_i}{b_i}} for all \m{i\in s},
then
\[\HasWeightNormLe{p}{\textstyle\sum_{i\in s}\kappa_i}{\textstyle\prod_{i\in s}F_i}{\textstyle\prod_{i\in s}b_i}.\]
\end{lemma}
\begin{proof}
By induction on~$s$. For the empty set the polynomial~$1$ is the polynomial of the coordinate vector of weight~$0$
with entry~$1$, of squared norm~$1$. For the step it suffices to treat two factors
\m{F=\operatorname{pol}_\kappa(\xi)} and \m{F'=\operatorname{pol}_{\kappa'}(\xi')}. The product \m{FF'} is semi-invariant of
weight \m{\kappa+\kappa'}, so by Input~1 it is \m{\operatorname{pol}_{\kappa+\kappa'}(\xi'')} for some~$\xi''$, and
by Input~2 \m{\|\xi''\|^2_{\kappa+\kappa'}\le\|\xi\|^2_\kappa\,\|\xi'\|^2_{\kappa'}}. The squared norms are \m{\ge0}, so the
bounds multiply.
\end{proof}
\leanline{\lean{HasWeightNormLeProd}}

\subsection*{Auxiliary pairs}
\begin{definition}[the finite groups of level places]\label{d:levels}
In Setting~\ref{set:bg}, let $\lpl$ be level places of~$E$, \m{\ix r} and~\m{\ix n} finite index sets and
\m{u\colon\ix r\hookrightarrow\ix n} an embedding.
\begin{enumerate}[label=\textup{(\roman*)},leftmargin=2.3em,itemsep=2pt]
\item \m{\ChosenGL{\lpl}{\ix r}} is the product of the groups \m{\GL_{\ix r}(k_w)} over the chosen places~$w$ of~$\lpl$.
\item \m{\firstBlockHom{\lpl}{u}{g}}, for \m{g=(g_w)\in\ChosenGL{\lpl}{\ix r}}, is the family in the product of the groups
\m{\GL_{\ix n}(k_w)} over all places of~$\lpl$ equal to \m{\operatorname{diag}(g_w,1)} at the chosen places, the block~$g_w$
on the indices \m{u(\ix r)}, and to~$1$ at the twist places; it is a homomorphism in~$g$.
\item \m{\auxLevel{\lpl}{u}} is the subgroup of the product of the groups \m{\GL_{\ix n}(k_w)} over all places of~$\lpl$
formed by the families that fix the basis vectors of index in \m{u(\ix r)} at the chosen places and whose rows of index in
\m{u(\ix r)} vanish outside the columns \m{u(\ix r)} at the twist places.
\end{enumerate}
\end{definition}
\leanline{\lean{ChosenGL}\sep \lean{firstBlockHom}\sep \lean{auxLevel}}

\begin{definition}[auxiliary pair]\label{d:aux}
In Setting~\ref{set:bg}, let \m{\ix r} and~\m{\ix n} be finite index sets. An \kw{auxiliary pair} for
\m{(E,\ix r,\ix n)}, an element~$O$ of \m{\AuxPair{E}{\ix r}{\ix n}}, consists of: level places of~$E$, with products
\m{\mathfrak{M}_O} and \m{\mathfrak{Q}_O}; a sign vector~$s_O$ on~$\ix r$ and a hermitian space~$W_O$ with index
set~$\ix r$ and signs~$s_O$; an embedding \m{e_O\colon\ix r\hookrightarrow\ix n}; a numbering
\m{\operatorname{ord}_O\colon\ix r\to\NN}; two weights $\kappa^{+}_O$ and~$\kappa^{-}_O$; and a character~$\mu_O$ with a weight
and a tame part~$\psi_O$ such that \m{\psi_{O,w}=1} at every finite place~$w$ that is not a place of the level
places; with \m{|\ix r|<|\ix n|}. The second space of the pair is~$I_{\ix n}$.
\end{definition}
\leanline{\lean{AuxPair}}

\begin{definition}[the objects of an auxiliary pair]\label{d:auxobj}
In Setting~\ref{set:bg}, let $O$ be an auxiliary pair for \m{(E,\ix r,\ix n)}.
\begin{enumerate}[label=\textup{(\arabic*)},leftmargin=2.3em,itemsep=2pt]
\item \emph{Level.} \m{\levelGroup{O}} is the level in \m{\KK{\ix n}} of the elements whose reductions at the places
of~$O$ form an element of \m{\auxLevel{O}{e_O}}. The character \m{\auxLevelChar{O}{g}} of \m{g\in\auxLevel{O}{e_O}} is the
product over the twist places~$w$ of \m{\psi_{O,w}^{-1}} at the determinant of the block of~$g_w$ with rows and
columns \m{e_O(\ix r)}, and \m{\levelChar{O}} is the character of \m{\levelGroup{O}} obtained by reduction.
\item \emph{Datum at the level.} For a family \m{X=(X_w)} of \m{\ix n\times\ix r} matrices over the residue fields at
the places of~$O$, \m{\resDetDatum{O}{X}} is the product over the twist places~$w$ of \m{\psi_{O,w}} at the determinant
of the rows \m{e_O(\ix r)} of~$X_w$ if \m{X_w=I_{e_O}} at every chosen place, and $0$ otherwise.
\item \emph{Weight and forms.} \m{\auxKappa{O}} is the weight and \m{\detExp{O}{\tau}} the integer given by
\[\auxKappa{O}\colon(\tau,i)\mapsto\kappa^{+}_{O,\tau}(i)+\kappa^{-}_{O,\tau}(|\ix n|-i),\qquad
\detExp{O}{\tau}=-\sum_{1\le i<|\ix n|}\kappa^{-}_{O,\tau}(i).\] \m{\FormV{O}} is the space of algebraic modular
forms on~$I_{\ix n}$ of weight \m{(\auxKappa{O},\detExp{O}{\cdot})} on which \m{\levelGroup{O}} acts through
\m{\levelChar{O}}. \m{\rankDiff{O}=|\ix n|-|\ix r|}.
\item \emph{The finite group.} For a complex representation~$\pi$ of \m{\UU(I_{\ix n})(\AAA_f)},
\m{\chosenGLRep{O}{\pi}} is the representation of \m{\ChosenGL{O}{\ix r}} on the vectors of~$\pi$ on which
\m{\levelGroup{O}} acts through \m{\levelChar{O}}, in which $g$ acts by \m{\prod_w\psi_{O,w}(\det g_w)}, the product
over the chosen places, times~$\pi$ of a lift of \m{\firstBlockHom{O}{e_O}{g}}. For the forms it is the representation
\m{\formAction{O}} of \m{\ChosenGL{O}{\ix r}} on \m{\FormV{O}}.
\item \emph{Polynomials.} \m{\tensorPoly{O}} is the linear map from the complex functions on pairs of standard
fillings, of weights $\kappa^{+}_O$ and~$\kappa^{-}_O$ with rows~$\ix n$, to polynomials that sends the indicator of a
pair \m{(l,l')} to the product of the polynomial of the indicator of~$l$, of weight~$\kappa^{+}_O$, written in the
positive letters and of the polynomial of the indicator of~$l'$, of weight~$\kappa^{-}_O$, written in the negative
letters: a variable \m{(\tau,i,h)} of the first is replaced by the sum of the variables \m{x_{\tau,i,c}} over the
\m{c\in\ix r} with \m{s_O(c)=+1} and \m{\operatorname{ord}_O(c)=h}, and likewise for the second with
\m{s_O(c)=-1}. The \kw{harmonic polynomial} of a coordinate vector~$\xi$ of weight \m{\auxKappa{O}} is
\m{\harmonicPoly{O}(\xi)}, where \m{\harmonicPoly{O}} is \m{\tensorPoly{O}} composed with the integral Cartan
embedding of the weight \m{\auxKappa{O}} into the tensor product of the weights $\kappa^{+}_O$ and~$\kappa^{-}_O$.
\item \emph{Kernel.} \m{\auxA{O}=\upsilon^{r_-}}, with $r_-$ the number of signs~$-1$ of~$s_O$, and
\m{\auxB{O}=\mu_O\,\upsilon^{\lfloor|\ix n|/2\rfloor}}. \m{\pairThetaKernel{O}} is the theta kernel of the pair
\m{(W_O,I_{\ix n})} with these two characters: the Weil representation \m{\omega_{\auxA{O},\auxB{O}}} with the theta
functional. \m{\auxDatum{O}} is the finite datum of the function
\m{x\mapsto\resDetDatum{O}{(x\bmod w)_w}} of the \m{\ix n\times\ix r} matrices modulo
\m{\mathfrak{M}_O^{2}\mathfrak{Q}_O}, and \m{\fockDatum{O}{\xi}} is the Schwartz function with archimedean part the Fock
vector of \m{\harmonicPoly{O}(\xi)} and finite part \m{\auxDatum{O}}.
\item \emph{Descent.} For \m{g\in\UU(W_O)(\AAA)} and \m{h\in\UU(I_{\ix n})(\AAA_f)}, \m{\conjKernel{O}{g}{h}} is the
linear form on coordinate vectors
\m{\xi\mapsto\overline{\theta\bigl(\omega_{\auxA{O},\auxB{O}}(g,h)\,\fockDatum{O}{\bar\xi}\bigr)}}. The \kw{descent} of a form
\m{\Xi\in\FormV{O}} is the function on \m{\UU(W_O)(\AAA)}
\[\descent{O}{\Xi}{g}=\clsum{h}{\levelGroup{O}}\conjKernel{O}{g}{h}\bigl(\Xi(h)\bigr).\]
\item \emph{Hecke test.} Let $w$ be a finite place with \m{w\ne\bar w} such that neither $w$ nor~$\bar w$ is a place
of~$O$. For \m{X\in\CC}, the operator \m{H_{O,w}(X)=\sum_{i=0}^{|\ix n|}(-X)^i q_w^{i(i-1)/2}T_{w,i}} acts on
\m{\FormV{O}}, the $T_{w,i}$ being the Hecke operators at~$w$. The \kw{test value} is
\m{\testValue{O}{w}=\bigl(\nu_{\auxA{O}}(w)\,q_w^{\rankDiff{O}-1}\bigr)^{-1}}, where \m{\nu_{\alpha}(w)} denotes the
Frobenius value at~$w$ of a character~$\alpha$ in the rank at hand. A form \m{y\in\FormV{O}} \kw{lifts with
parameters~$z$}, where $z$ assigns to each finite place a sequence of complex numbers, if at every finite
place~$w$ as above $y$ is an eigenvector of all \m{H_{O,w}(X)}, with eigenvalue
\[\prod_{0\le l<\rankDiff{O}}\bigl(1-q_w^{l}\,\nu_{\auxA{O}}(w)X\bigr)\cdot\sum_{i=0}^{|\ix r|}\Bigl(-\frac{q_w^{\rankDiff{O}}\,\nu_{\auxA{O}}(w)}{\nu_{\auxB{O}}(w)}\,X\Bigr)^{i}q_w^{i(i-1)/2}z(w)_i .\]
\end{enumerate}
\end{definition}
\leanline{\lean{levelGroup}\sep \lean{auxLevelChar}\sep \lean{levelChar}\sep \lean{resDetDatum}\sep \lean{auxKappa}\sep
\lean{detExp}\sep \lean{FormV}\sep \lean{rankDiff}\sep \lean{chosenGLRep}\sep \lean{formAction}\sep \lean{tensorPoly}\sep
\lean{harmonicPoly}\sep \lean{auxA}\sep \lean{auxB}\sep \lean{pairThetaKernel}\sep \lean{auxDatum}\sep \lean{fockDatum}\sep
\lean{conjKernel}\sep \lean{descent}\sep \lean{heckePolyOnForms}\sep \lean{testValue}\sep \lean{Lifts}}

\begin{remark}[two small facts]\label{r:auxfacts}
(i) For coordinate vectors $a$ and~$b$ of weights $\kappa^{+}_O$ and~$\kappa^{-}_O$, the map \m{\tensorPoly{O}} sends the
function \m{(l,l')\mapsto a_l\,b_{l'}} to the product of the polynomial of~$a$ in the positive letters and the
polynomial of~$b$ in the negative letters; this is the definition, extended bilinearly. (ii) If a form~$y$ in
\m{\FormV{O}} is, at a place~$w$ as in Definition~\ref{d:auxobj}\,(8), an eigenvector of the Hecke polynomial of the
right regular representation with eigenvalue \m{p(X)}, then \m{H_{O,w}(X)y=p(X)y}: the operator on the forms is the
restriction of that Hecke polynomial.
\end{remark}
\leanline{\lean{tensorPolyVec}\sep \lean{heckePolySmul}}

The auxiliary pairs of the proof are the pairs \m{(W_j,V_j)} of the single lifts: the adjoint lift of the note
\emph{The volume of the chain of theta lifts as a product of lift volumes} is the adjoint of the descent of
Definition~\ref{d:auxobj}\,(7), and the Hecke test enters the notion of a fit family of twists.

\subsection*{Residue rings and the divided splitting}
\begin{definition}[moduli of level places]\label{d:moduli}
In Setting~\ref{set:bg}, let $\lpl$ be level places of~$E$. \m{\chosenModulusSq{\lpl}} is the modulus
\m{\mathfrak{M}_\lpl^{2}}; for \m{x\in\OEr/\mathfrak{M}_\lpl^{2}\mathfrak{Q}_\lpl}, \m{\reduceToRzero{x}} is its image under the
reduction to \m{\OEr/\mathfrak{M}_\lpl^{2}}; and \m{\unitModulus{E}} is the modulus whose ideal is~$\OEr$.
\end{definition}
\leanline{\lean{chosenModulusSq}\sep \lean{reduceToRzero}\sep \lean{unitModulus}}

\begin{definition}[generator and divided splitting]\label{d:div}
In Setting~\ref{set:bg}, let $O$ be an auxiliary pair for \m{(E,\ix r,\ix n)}.
\begin{enumerate}[label=\textup{(\roman*)},leftmargin=2.3em,itemsep=2pt]
\item \m{\genIdele{O}} is the chosen generator of the modulus \m{\mathfrak{M}_O}: an integral idele, equal to~$1$ at
the places that are not chosen places of~$O$; \m{\genIntAdele{O}} is the same element, regarded in~$\Ohat$.
\item Let \m{\ix N}, \m{\ix m}, \m{\ix a} be finite index sets and \m{z\colon\ix r\hookrightarrow\ix m} an embedding.
A \kw{divided splitting}, an element~$\mathcal{V}$ of \m{\DividedSplitting{\ix N}{\ix m}{\ix a}{O}{z}}, consists of a
bijection \m{\varepsilon\colon\ix N\sqcup\ix m\to\ix n} with \m{\varepsilon(z(i))=e_O(i)} for \m{i\in\ix r}, of two
embeddings \m{\ix r\hookrightarrow\ix N} and \m{\ix a\hookrightarrow\ix N}, the \kw{fixed} and the \kw{divided}
embedding, and of the conditions that \m{|\ix N|\le|\ix n|} and $\kappa^{+}_O$ has no column of height~$h$ with
\m{|\ix N|\le h<|\ix n|}, and the same for $\ix m$ and~$\kappa^{-}_O$.
\item For a divided splitting~$\mathcal{V}$, \m{\genResidue{\mathcal{V}}} is the image of \m{\genIntAdele{O}} in
\m{\OEr/\mathfrak{M}_O^{2}}; it depends on~$O$ only.
\end{enumerate}
\end{definition}
\leanline{\lean{genIdele}\sep \lean{genIntAdele}\sep \lean{DividedSplitting}\sep \lean{genResidue}}

\begin{lemma}[reduction modulo a power of the chosen primes]\label{l:red}
In Setting~\ref{set:bg}, let $O$ be an auxiliary pair for \m{(E,\ix r,\ix n)}, \m{c\ge0} an integer and
\m{x\in\Ohat}. Then $x$ reduces to~$0$ modulo \m{\mathfrak{M}_O^{c}} if and only if
\m{\genIdele{O}^{-c}\,x\in\Ohat}.
\end{lemma}
\begin{proof}
The kernel of the reduction modulo \m{\mathfrak{M}_O^{c}} consists of the integral adeles of valuation at
least~$c$ at every chosen place. The idele \m{\genIdele{O}} has valuation~$1$ at the chosen places, which are prime
factors of multiplicity one of~\m{\mathfrak{M}_O}, and is~$1$ elsewhere. So \m{\genIdele{O}^{-c}x} is integral at every place
exactly when $x$ has valuation at least~$c$ at the chosen places.
\end{proof}
\leanline{\lean{reductionZero}}

\begin{lemma}[multiplying by the generator tests congruence modulo the chosen primes]\label{l:genres}
In Setting~\ref{set:bg}, let $O$ be an auxiliary pair for \m{(E,\ix r,\ix n)}, \m{\ix N}, \m{\ix m} finite index
sets, \m{\ix a} an index set, \m{z\colon\ix r\hookrightarrow\ix m} and $\mathcal{V}$ a divided splitting in
\m{\DividedSplitting{\ix N}{\ix m}{\ix a}{O}{z}}. Let \m{\ix k}, \m{\ix c} be index sets, $Z$ a \m{\ix k\times\ix c}
matrix over \m{\OEr/\mathfrak{M}_O^{2}\mathfrak{Q}_O} and \m{u\colon\ix c\to\ix k} a map, with $I_u$ the matrix whose
column~$i$ is the standard basis vector of index~\m{u(i)}. Then
\[\genResidue{\mathcal{V}}\cdot\reduceToRzero{Z}=\genResidue{\mathcal{V}}\cdot I_u\quad\Longleftrightarrow\quad
Z\equiv I_u\ \text{modulo every chosen place of }O,\]
where \m{\reduceToRzero{Z}} is $Z$ reduced entry by entry.
\end{lemma}
\begin{proof}
Write each entry of \m{Z-I_u} as the reduction of an integral adele, and let $X$ be the matrix of these adeles.
The left side says that \m{\genIntAdele{O}X} reduces to~$0$ modulo \m{\mathfrak{M}_O^{2}}. By Lemma~\ref{l:red},
for \m{c=2} and for \m{c=1}, this holds if and only if $X$ reduces to~$0$ modulo~\m{\mathfrak{M}_O}. Since
\m{\mathfrak{M}_O} is a product of distinct primes, two matrices over \m{\OEr/\mathfrak{M}_O} are equal if and only if
they are equal modulo each of these primes (Input~3). This is the right side.
\end{proof}
\leanline{\lean{genResidueSmul}}

\begin{remark}[bookkeeping lemmas of the file]\label{r:book}
Five further lemmas of the same file are used in the other notes of this unit and need no argument. For level
places~$\lpl$: the ideal \m{\mathfrak{M}_\lpl^{2}\mathfrak{Q}_\lpl} is contained in \m{\mathfrak{M}_\lpl^{2}}; the primes dividing
\m{\mathfrak{M}_\lpl^{c}}, \m{c\ge1}, are the chosen places; the chosen lift to \m{\KK{\ix n}} of an invertible matrix
over \m{\OEr/\mathfrak{M}_\lpl^{2}} has component~$1$ outside the chosen places and their conjugates, and an element with
this property has component~$1$ at every twist place; and reduction modulo a ring homomorphism of a block diagonal
element of \m{\KK{\ix n}}, along a bijection \m{\ix N\sqcup\ix m\to\ix n}, is the block diagonal matrix of the
reductions. For an auxiliary pair~$O$: \m{\genIdele{O}} has component~$1$ at every finite place that is not a chosen
place.
\end{remark}
\leanline{\lean{levelModulusIdealLe}\sep \lean{placesLevelModulus}\sep \lean{liftModSupported}\sep
\lean{atPlaceTwist}\sep \lean{mapBlock}\sep \lean{genIdeleApply}}

\subsection*{Stacked determinants}
\begin{definition}[stacked determinants, substitutions, coordinates]\label{d:stacked}
Let $\Plc$ be a finite set.
\begin{enumerate}[label=\textup{(\roman*)},leftmargin=2.3em,itemsep=2pt]
\item Let \m{\ix l} and~\m{\ix n'} be finite index sets, \m{\ix n} an index set, \m{\tau\in\Plc}, $\eta$ a complex
\m{\ix l\times\ix n} matrix and \m{j\colon\ix l\sqcup\ix n'\to\ix n} a map. The \kw{stacked determinant}
\m{\stackedDet{\tau}{\eta}{j}} is the determinant of the square matrix, with rows and columns indexed by
\m{\ix l\sqcup\ix n'}, whose entry at \m{(i,c)} is \m{\eta_{i,j(c)}} for \m{i\in\ix l} and the variable
\m{x_{\tau,i,j(c)}} for \m{i\in\ix n'}: the columns~$j$ of~$\eta$ stacked over the matrix~$x_\tau$ of variables. It
is a complex polynomial in the variables \m{x_{\tau,i,c}}, \m{i\in\ix n'}, \m{c\in\ix n}.
\item Let \m{\ix a}, \m{\ix N}, \m{\ix L} be finite index sets, \m{e\colon\ix a\sqcup\ix N\to\ix L} a bijection,
\m{(\iota_\tau)_{\tau\in\Plc}} a family of finite sets and \m{M=(M_{\tau k})} a family of complex \m{\ix a\times\ix L}
matrices, \m{k\in\iota_\tau}. The \kw{product of stacked determinants} is
\[\sdp{e}{M}=\prod_{\tau\in\Plc}\ \prod_{k\in\iota_\tau}\stackedDet{\tau}{M_{\tau k}}{e},\]
a polynomial in the variables \m{x_{\tau,i,c}}, \m{i\in\ix N}, \m{c\in\ix L}.
\item For families \m{c=(c_\tau)} of complex \m{\ix N\times\ix N} matrices and \m{g=(g_\tau)} of complex
\m{\ix L\times\ix L} matrices, \m{\sub{c}{g}} is the substitution \m{x_\tau\mapsto c_\tau\transp x_\tau g_\tau\transp} on
the polynomials in these variables: \m{x_{\tau,a,b}\mapsto\sum_{i,j}c_\tau(i,a)\,g_\tau(b,j)\,x_{\tau,i,j}}.
\item Let \m{p\ge0} be an integer. For a polynomial~$F$ in the variables \m{x_{\tau,a,b}}, with \m{a\in\ix N} and $b$
in the disjoint union of \m{\{0,\dots,p-1\}} and~$\ix L$, \m{\infc(F)} is $F$ regarded as a polynomial in the
variables of the first $p$ columns, with coefficients in the polynomials in the variables of the columns~$\ix L$.
\item For a weight~$\kappa$, such a polynomial~$F$ and a standard filling~$l$ of weight~$\kappa$ with
rows~$\ix N$, the \kw{$l$-th coordinate of~$F$ in the first columns} is
\[\firstColsCoord{\kappa}{F}{l}=h_\kappa^{-1}\cdot\ip{\infc(F)}{b_l},\]
the Fock pairing taken in the variables of the first $p$ columns; it is a polynomial in the variables of the
columns~$\ix L$.
\end{enumerate}
\end{definition}
\leanline{\lean{stackedDet}\sep \lean{stackedDetProd}\sep \lean{twoSidedSubst}\sep \lean{inFirstCols}\sep
\lean{firstColsCoord}}

\begin{lemma}[a determinant with fixed columns]\label{l:detcols}
Let \m{\ix a}, \m{\ix N}, \m{\ix L} be finite index sets, \m{e\colon\ix L\to\ix a\sqcup\ix N} a bijection and $Z$ a
complex \m{\ix L\times\ix N} matrix. Then there are a complex \m{\ix a\times\ix L} matrix~$Y$ and \m{\gamma\in\CC} such
that for every complex \m{\ix L\times\ix a} matrix~$X$
\[\det\bigl[X\mid Z\bigr]_e=\gamma\cdot\det\bigl((YX)\transp\bigr),\qquad\text{and}\qquad|\gamma|^2\cdot|\det(YY\herm)|\le|\det(Z\herm Z)| ,\]
where \m{[X\mid Z]_e} is the square matrix with columns $X$ and~$Z$, its columns indexed by~$\ix L$ through~$e$.
\end{lemma}
\begin{proof}
By Input~4 there is an invertible \m{\ix L\times\ix L} matrix whose columns, indexed through~$e$, split as
\m{[U_1\mid U_2]} with \m{Z\herm U_1=0}, that is \m{U_1\herm Z=0}. For matrices $A$, $X$ of size \m{\ix L\times\ix a}
and $B$ of size \m{\ix L\times\ix N},
\[\overline{\det[A\mid B]_e}\cdot\det[X\mid Z]_e=\det\begin{pmatrix}A\herm X&A\herm Z\\ B\herm X&B\herm Z\end{pmatrix}.\]
For \m{A=U_1} and \m{B=U_2} the upper right block vanishes, so
\[\overline{\det[U_1\mid U_2]_e}\cdot\det[X\mid Z]_e=\det(U_1\herm X)\det(U_2\herm Z):\]
the first assertion holds with \m{Y=U_1\herm} and \m{\gamma=\det(U_2\herm Z)/\overline{\det[U_1\mid U_2]_e}}. For \m{A=X=U_1} and \m{B=Z} the same identity gives
\[|\det[U_1\mid Z]_e|^2=\det(YY\herm)\det(Z\herm Z),\]
and the left side is \m{|\gamma|^2\det(YY\herm)^2}. If \m{\det(YY\herm)\ne0}, divide by it; otherwise the inequality is trivial.
\end{proof}
\leanline{\lean{detFromCols}}

\begin{lemma}[the substitution law]\label{l:subst}
Let $\Plc$, \m{\ix a}, \m{\ix N}, \m{\ix L} be finite, \m{e\colon\ix a\sqcup\ix N\to\ix L} a bijection,
\m{(\iota_\tau)} a family of finite sets, \m{M=(M_{\tau k})} a family of complex \m{\ix a\times\ix L} matrices and
\m{c=(c_\tau)}, \m{g=(g_\tau)} families of complex \m{\ix N\times\ix N} and \m{\ix L\times\ix L} matrices. Then
\[\sub{c}{g}\Bigl(\sdp{e}{(M_{\tau k}\,g_\tau\transp)_{\tau,k}}\Bigr)=\prod_{\tau\in\Plc}\bigl(\det c_\tau\cdot\det g_\tau\bigr)^{|\iota_\tau|}\cdot\sdp{e}{M}.\]
\end{lemma}
\begin{proof}
The substitution replaces the matrix \m{M_{\tau k}g_\tau\transp} stacked over~$x_\tau$ by \m{M_{\tau k}g_\tau\transp}
stacked over \m{c_\tau\transp x_\tau g_\tau\transp}, which is \m{\operatorname{diag}(1,c_\tau\transp)} times the matrix
\m{M_{\tau k}} stacked over~$x_\tau$, times~\m{g_\tau\transp}. Determinants are multiplicative, and the columns are
ordered by the same bijection~$e$ on both sides. Each of the \m{|\iota_\tau|} factors at~$\tau$ contributes
\m{\det c_\tau\det g_\tau}.
\end{proof}
\leanline{\lean{substStackedDetProd}}

\begin{lemma}[coordinates in the first columns]\label{l:coord}
Let $\Plc$, \m{\ix N}, \m{\ix L} be finite, \m{p\ge0} an integer, \m{c=(c_\tau)} and \m{g=(g_\tau)} families of complex
\m{\ix N\times\ix N} and \m{\ix L\times\ix L} matrices, $F$ and~$F'$ polynomials as in
Definition~\ref{d:stacked}\,(iv) and \m{s\in\CC}. Suppose
\m{\sub{c}{\operatorname{diag}(1,g)}(F')=s\cdot F}, where \m{\operatorname{diag}(1,g)} is the family of the block
matrices with the identity on the first $p$ columns and~$g_\tau$ on~$\ix L$. Then for every weight~$\kappa$ and
every standard filling~$l$ of weight~$\kappa$ with rows~$\ix N$
\[\sub{c}{g}\Bigl(\sum_{l'}S_\kappa(c)_{l,l'}\cdot\firstColsCoord{\kappa}{F'}{l'}\Bigr)=s\cdot\firstColsCoord{\kappa}{F}{l}.\]
\end{lemma}
\begin{proof}
Apply \m{\infc} to the hypothesis. On the variables of the first $p$ columns the substitution
\m{\sub{c}{\operatorname{diag}(1,g)}} is the row substitution \m{y\mapsto c\transp y}, and on the coefficients it is
\m{\sub{c}{g}}. So \m{s\cdot\infc(F)} is obtained from \m{\infc(F')} by applying \m{\sub{c}{g}} to the coefficients
and the row substitution by~$c\transp$ to the variables. Pair with~$b_l$. A ring homomorphism applied to the
coefficients commutes with the Fock pairing, and the row substitution by~$c\transp$ on one side is the row
substitution by~$c$ on the other (Input~5): the pairing becomes \m{\sub{c}{g}} of \m{\ip{\infc(F')}{b_l(c\,y)}}. By the
definition of the Schur matrix \m{b_l(cy)=\sum_{l'}S_\kappa(c)_{l,l'}b_{l'}(y)}. Dividing by~$h_\kappa$ gives the claim.
\end{proof}
\leanline{\lean{substFirstColsCoord}}

\subsection*{The carrier}
\begin{definition}[the carrier]\label{d:carrier}
In Setting~\ref{set:bg}, let \m{m\ge0} be an integer.
\begin{enumerate}[label=\textup{(\roman*)},leftmargin=2.3em,itemsep=2pt]
\item The \kw{carrier} \m{\carrier{E}{m}} is the theta lift of the constant function~$1$ from the definite group
\m{\UU(I_m^{-})} to the unitary group of signature \m{(2m,m)}: for \m{x\in\UU(I_{2m,m})(\AAA)},
\[\carrier{E}{m}(x)=\int_{[\UU(I_m^{-})]}\tk{I_m^{-}}{I_{2m,m}}{\upsilon^{m}}{1}{\unitModulus{E}}{1}{\textstyle\prod_{\tau}\det(y^{-}_\tau)^{2m}}(g,x)\,dg,\]
where the product runs over the infinite places~$\tau$ and \m{y^{-}_\tau} is the \m{m\times m} matrix of the
variables \m{x_{\tau,i,c}} with $i$ in the negative block of the index set of \m{I_{2m,m}}.
\item For a character~$c$, an index set~$\ix n'$ and \m{h\in\UU(I_{\ix n'})(\AAA_f)},
\[\carrierXi{m}{c}{h}=\clsum{u}{\KK{m}}\tk{I_m}{I_{\ix n'}}{c}{1}{\unitModulus{E}}{1}{1}(u,h),\]
the class sum over \m{\UU(I_m)} of the theta kernel with constant datum and polynomial~$1$.
\end{enumerate}
\end{definition}
\leanline{\lean{carrierForm}\sep \lean{carrierXi}}

\subsection*{The archimedean matching}
\begin{definition}[block signs, the positive block, the matched polynomial]\label{d:real}
Let $E$ be a number field and \m{\ix n,\ix m,\ix a,\ix b,\ix p,\ix l} finite index sets.
\begin{enumerate}[label=\textup{(\roman*)},leftmargin=2.3em,itemsep=2pt]
\item For a bijection \m{f\colon\ix p\sqcup\ix l\to\ix b} and \m{i\in\ix b}, the sign \m{\blockSigns{f}{i}} is $+1$ if
\m{i\in f(\ix p)} and $-1$ if \m{i\in f(\ix l)}.
\item For a bijection \m{j\colon\ix m\sqcup\ix l\to\ix n} and \m{k\in G_\infty(1)}, the definite group on~$\ix m$,
the element \m{\posBlockEmb{j}{k}} of \m{G_\infty(\blockSigns{j}{\cdot})} is, at each infinite place~$\tau$, the matrix
with block~$k_\tau$ on \m{j(\ix m)} and the identity on \m{j(\ix l)}; it is a homomorphism in~$k$.
\item For bijections $f$ and~$j$ as in (i) and~(ii), a family \m{(\iota_\tau)} of finite sets indexed by the infinite
places and a family \m{\eta=(\eta_{\tau k})} of complex \m{\ix p\times\ix m} matrices, \m{k\in\iota_\tau}, the
\kw{matched polynomial} is
\[\mpoly{f}{j}{\eta}=\prod_{\tau}\det\bigl(y^{--}_\tau\bigr)^{|\iota_\tau|+|\ix n|-2|\ix l|}\cdot\prod_{\tau}\prod_{k\in\iota_\tau}\det\bigl(\eta_{\tau k}\,y^{++}_\tau\bigr),\]
a polynomial in the variables \m{x_{\tau,i,c}}, \m{i\in\ix n}, \m{c\in\ix b}, where \m{y^{--}_\tau} is the block of the
variables with \m{i\in j(\ix l)}, \m{c\in f(\ix l)} and \m{y^{++}_\tau} the block with \m{i\in j(\ix m)},
\m{c\in f(\ix p)}. The exponent is a truncated difference of natural numbers.
\item For a bijection~$j$ as in~(ii), \m{\posFockVector{j}} is the linear map from the polynomials in the variables
\m{x_{\tau,i,c}}, \m{i\in\ix a}, \m{c\in\ix m}, to the Fock space of the pair of signs \m{\blockSigns{j}{\cdot}} on~$\ix n$
and of the definite space on~$\ix a$: replace the column index~$c$ by \m{j(c)} and take the Fock vector.
\end{enumerate}
\end{definition}
\leanline{\lean{blockSigns}\sep \lean{posBlockEmb}\sep \lean{matchedPoly}\sep \lean{posFockVector}}

\begin{remark}\label{r:pos}
Two computations about these objects are used without comment below: \m{\posBlockEmb{j}{k}} is block diagonal for the
signs \m{\blockSigns{j}{\cdot}}, its inverse is \m{\posBlockEmb{j}{k^{-1}}}, its minor on the positive block is
\m{\det k_\tau} and its minor on the negative block is~$1$; and the number of \m{i\in\ix b} with
\m{\blockSigns{f}{i}=-1} is~\m{|\ix l|}.
\end{remark}
\leanline{\lean{posBlockEmbSpec}\sep \lean{cardNegBlockSigns}}

\inputbox{\textbf{Inputs.} The results below are taken for granted. Inputs 1 to~8 are proved in the background
library of the artifact; the catalogue of the handout records, for each, whether a reference with its place has
been confirmed. Input~9 belongs to the unit \emph{Exchange} of this chapter.

\vspace{4pt}\textbf{Input 1} (the image of \m{\operatorname{pol}_\kappa}).\enspace Over~$\CC$, the image of
\m{\xi\mapsto\operatorname{pol}_\kappa(\xi)} is the space of all polynomials that are semi-invariant of weight~$\kappa$;
the product of semi-invariants of weights $\kappa$ and~$\kappa'$ is semi-invariant of weight \m{\kappa+\kappa'}.
(Part of the standard basis theorem for Schur modules; reference to be confirmed.)

\vspace{3pt}\textbf{Input 2} (norms of products).\enspace If
\m{\operatorname{pol}_{\kappa+\kappa'}(\xi'')=\operatorname{pol}_\kappa(\xi)\operatorname{pol}_{\kappa'}(\xi')}, then
\m{\|\xi''\|^2_{\kappa+\kappa'}\le\|\xi\|^2_\kappa\,\|\xi'\|^2_{\kappa'}}. (Reference to be confirmed.)

\vspace{3pt}\textbf{Input 3} (the Chinese remainder theorem for matrices).\enspace If the ideal of a
modulus is a product of distinct primes, two matrices over its residue ring are equal if and only if their
reductions to the residue fields of these primes are equal.

\vspace{3pt}\textbf{Input 4} (clearing columns; the Hermite normal form over a field).\enspace For a complex
\m{\ix N\times\ix L} matrix~$M$ and an embedding \m{\ix N\hookrightarrow\ix L} there is an invertible
\m{\ix L\times\ix L} matrix~$U$ such that the columns of~$MU$ with index outside the image of the embedding vanish.
(The library proves it for valuation rings; a place for the reference is not confirmed.)

\vspace{3pt}\textbf{Input 5} (the Fock pairing).\enspace (a) For matrices~$z_\tau$, the row substitutions by $z$ and
by~$z\transp$ are adjoint: \m{\ip{p(zy)}{q}=\ip{p}{q(z\transp y)}}. (b) A ring homomorphism applied to the coefficients of
two polynomials commutes with their Fock pairing. (c) For a non-zero complex polynomial~$q$, with $\bar q$ the
polynomial with conjugated coefficients, \m{\ip{\bar q}{q}>0}.

\vspace{3pt}\textbf{Input 6} (the real Weil representation on Fock vectors).\enspace Let $s$, $s'$ be sign vectors of
signatures \m{(p_0,q_0)} and \m{(p_0',q_0')}, and let \m{x\in G_\infty(s)}, \m{y\in G_\infty(s')} be block diagonal, with
blocks \m{x^{\pm}_\tau}, \m{y^{\pm}_\tau}. Then for every polynomial~$P$
\[\omega^{\infty}_{\gamma,\gamma'}(x,y)\,v_P=\prod_\tau(\det y^{+}_\tau)^{\gamma_\tau-q_0}(\det y^{-}_\tau)^{\gamma_\tau-p_0}(\det x^{+}_\tau)^{q_0'-\gamma'_\tau}(\det x^{-}_\tau)^{p_0'-\gamma'_\tau}\cdot v_{P'},\]
where $P'$ is obtained from~$P$ by the substitution \m{X\mapsto y\transp X\,(x^{-1})\transp} on the blocks of variables
whose row and column have the same sign and \m{X\mapsto y^{-1}X\,x} on the other blocks.

\vspace{3pt}\textbf{Input 7} (polarisation).\enspace Let $\mathcal{S}$ be a commutative $\CC$-algebra, $M$ a linear
form on~$\mathcal{S}$, \m{Y=(Y_\tau)} matrices over~$\mathcal{S}$ of size \m{\ix n\times\ix a}, $q$ a complex polynomial
in the entries of a generic family \m{(X'_\tau)} of \m{\ix n\times\ix a} matrices, $K$ a complex polynomial in the
entries of a generic family \m{(X_\tau)} of \m{\ix n\times\ix n} matrices, and \m{r\colon\ix a\to\ix n} a map. If
\m{M\bigl(q(u\transp Y)\bigr)=K(u)} for every family \m{u=(u_\tau)} of unitary \m{\ix n\times\ix n} matrices, then for
every family \m{\eta=(\eta_{\tau k})}, \m{k\in\iota_\tau}, of complex \m{\ix a\times\ix n} matrices
\[\Bigl\langle q,\ \prod_\tau\det\bigl(X'^{\,r}_\tau\bigr)^{|\iota_\tau|}\Bigr\rangle\cdot M\Bigl(\prod_{\tau,k}\det(\eta_{\tau k}Y_\tau)\Bigr)
=\Bigl\langle K,\ \prod_{\tau,k}\det(\eta_{\tau k}X^{r}_\tau)\Bigr\rangle,\]
where \m{X'^{\,r}_\tau} is the \m{\ix a\times\ix a} block of~$X'_\tau$ in the rows~$r$ and \m{X^{r}_\tau} the
\m{\ix n\times\ix a} block of~$X_\tau$ in the columns~$r$. (Reference to be confirmed.)}

\inputbox{\textbf{Inputs, continued.}

\vspace{4pt}\textbf{Input 8} (calculus of theta kernels).\enspace In Setting~\ref{set:bg}:
(a) \emph{decay}: if \m{\mathfrak{m}\ne(1)} and \m{\Psi(y)\ne0} only for matrices~$y$ with a left inverse, then for
every polynomial~$P$ and every~$h$ the function \m{g\mapsto\tk{S}{T}{\alpha}{\beta}{\mathfrak{m}}{\Psi}{P}(g,h)} decreases
rapidly on \m{\UU(S)(\AAA)}; this is Lemma~\ref{l:rapid} below, a direct consequence of the decay estimate of the
library and of the finiteness of the moments of~$v_P$.
(b) \emph{invariance}: the theta functional is invariant under the rational points of both groups.
(c) \emph{twisting}: for characters $\mu$, $\mu'$,
\m{\tk{S}{T}{\alpha\mu}{\beta\mu'}{\mathfrak{m}}{\Psi}{P}(g,h)=\mu(\det h)\,\mu'(\det g)^{-1}\,\tk{S}{T}{\alpha}{\beta}{\mathfrak{m}}{\Psi}{P}(g,h)}.
(d) \emph{translation by integral points}: right translation of $g$ or of~$h$ by an element~$k$ of the stabiliser
of the standard lattice can be moved into the datum: it replaces \m{\Psi} by a character value at \m{\det k} times
the translate of~$\Psi$ by the reduction of~$k$ modulo~$\mathfrak{m}$; for a tame character the value is that of
its character~$\chi$ at the determinant of the reduction.
(e) \emph{seesaw}: for an orthogonal sum in the second place, the kernel of a product of data on the two blocks of
rows and of a product of polynomials is the product of the kernels,
\m{\tk{S}{T\oplus V}{\alpha}{\beta\beta'}{\mathfrak{m}}{\Psi\Psi'}{PP'}(g,(h,h'))
=\tk{S}{T}{\alpha}{\beta}{\mathfrak{m}}{\Psi}{P}(g,h)\cdot\tk{S}{V}{\alpha}{\beta'}{\mathfrak{m}}{\Psi'}{P'}(g,h')}, and likewise
for a sum in the first place.
(f) \emph{change of modulus}: if \m{\mathfrak{m}\subseteq\mathfrak{n}}, the kernel modulo~$\mathfrak{m}$ of a
function that factors through the matrices modulo~$\mathfrak{n}$ is the kernel modulo~$\mathfrak{n}$ of that
function.
(g) \emph{negation}: there is a number of absolute value~$1$, depending on $S$ and~$T$ only, by which the kernel of
\m{(S^{-},T)}, with the inverse characters times powers of~$\upsilon$, the conjugated polynomial and the conjugated
datum composed with multiplication by an element of norm~$-1$, differs from the complex conjugate of the kernel of
\m{(S,T)}.
(h) \emph{definite groups}: for a function on \m{[\UU(I_{\ix a})]} that is invariant under \m{G_\infty} and under
\m{\KK{\ix a}}, the integral is the class sum at level~\m{\KK{\ix a}}; a class sum at a smaller level~$\Gamma$ is
the class sum at level \m{\KK{\ix a}} of the sum of the translates by representatives of \m{\KK{\ix a}/\Gamma}.

\vspace{3pt}\textbf{Input 9} (unit \emph{Exchange}; the exchange theorem and block partners).\enspace
(a) Let $W$, $V$, $V'$ be hermitian spaces with index sets \m{\ix n}, \m{\ix a}, \m{\ix b} and
\m{\sigma\colon\ix b\sqcup\ix a\to\ix n} a bijection; suppose that the number~$q_W$ of negative signs of~$W$ is the
sum of those of $V$ and~$V'$, that \m{|\ix a|=|\ix b|+2d}, that \m{2q_W+2\le|\ix n|}, and let $\mu$ be a weight with
\m{|\ix n|\le2\mu_\tau+2q_W} for all~$\tau$ and $v$ a vector of the Fock space of \m{(V',W)} with finite moments.
Then there are a real number~$C_0$ and a vector~$x_0$ of the Fock space of \m{(W,V)} with finite moments such that
for every modulus~$\mathfrak{m}$, character~$\nu$ and character~$\chi$ of \m{\OEr/\mathfrak{m}} with $\nu$ tame for
\m{(\mathfrak{m},\mu,\chi)} there is \m{c_0\ne0} with \m{|c_0|^2=C_0\cdot\LP{\mathfrak{m}}{\nu}{d}} and, for all partners
$\dL$, $\Psi$ for \m{(\sigma,\chi)} and all \m{g\in\UU(W)(\AAA)}, with
\m{\alpha_0=\nu^{-1}\upsilon^{\lfloor|\ix n|/2\rfloor}},
\[\int_{[\UU(V)]}\theta\bigl(\omega_{\alpha_0,1}(g,h)(x_0\otimes\dL)\bigr)\,dh
=c_0\cdot\int_{[\UU(V')]}\theta\bigl(\omega_{\alpha_0,\alpha_0}(h,g)(v\otimes\Psi)\bigr)\,dh .\]
(b) For a multiplicative character~$\chi$ of a commutative ring and any function~$f$, the functions
\m{\dL(x_1\mid x_2)=\chi^{-1}(\det x_2)\,f(x_2^{-1}x_1)} and
\m{\Psi\bigl(\begin{smallmatrix}y_1\\y_2\end{smallmatrix}\bigr)=\chi^{-1}(\det y_1)\,f(-y_2y_1^{-1})}, with $x_2$ and~$y_1$
square, are partners for the identity bijection and~$\chi$; and partners $\dL$, $\Psi$ for \m{(\sigma,\chi)} satisfy
\m{\chi(\det k)\,\dL(kx)=\dL(x)} for invertible~$k$.}

\sect{Proofs}{s:proofs}
\subsection*{The archimedean step}
\begin{lemma}[equivariance of the product of stacked determinants]\label{l:realweil}
Let $E$ be a number field, \m{\ix n,\ix m,\ix a,\ix p,\ix l} finite index sets, \m{j\colon\ix m\sqcup\ix l\to\ix n} and
\m{e\colon\ix p\sqcup\ix a\to\ix m} bijections, \m{(\iota_\tau)} a family of finite sets indexed by the infinite
places, $\gamma$ the weight with \m{\gamma_\tau=|\ix l|-|\iota_\tau|}, and \m{\eta=(\eta_{\tau k})} a family of complex
\m{\ix p\times\ix m} matrices. For all \m{k\in G_\infty(1)}, the definite group on~$\ix m$, and all~$u$ in the
definite group on~$\ix a$,
\[\omega^{\infty}_{\gamma,0}\bigl(\posBlockEmb{j}{k},u\bigr)\ \posFockVector{j}\Bigl(\sdp{e}{(\eta_{\tau k'}\,(k_\tau^{-1})\transp)_{\tau,k'}}\Bigr)
=\prod_\tau(\det k_\tau)^{-|\iota_\tau|}\cdot\posFockVector{j}\bigl(\sdp{e}{\eta}\bigr).\]
\end{lemma}
\begin{proof}
By Input~6 for the pair of signs \m{\blockSigns{j}{\cdot}}, of signature \m{(|\ix m|,|\ix l|)}, and of the definite signs
on~$\ix a$: the element \m{(\posBlockEmb{j}{k},u)} multiplies a Fock vector by
\m{\prod_\tau(\det u_\tau)^{\gamma_\tau-|\ix l|}=\prod_\tau(\det u_\tau)^{-|\iota_\tau|}} and substitutes
\m{x_\tau\mapsto u_\tau\transp x_\tau(k_\tau^{-1})\transp} in the variables of the positive columns, the only ones that
occur. By Lemma~\ref{l:subst}, with \m{c=u} and \m{g=k^{-1}}, this substitution multiplies the polynomial on the left
by \m{\prod_\tau(\det u_\tau\det k_\tau^{-1})^{|\iota_\tau|}} and turns it into \m{\sdp{e}{\eta}}. The powers of
\m{\det u_\tau} cancel.
\end{proof}
\leanline{\lean{realWeilStackedDetProd}}

\begin{proposition}[the archimedean matching]\label{p:real}
Let $E$ be a number field, \m{\ix n,\ix m,\ix a,\ix b,\ix p,\ix l} finite index sets,
\m{f\colon\ix p\sqcup\ix l\to\ix b}, \m{j\colon\ix m\sqcup\ix l\to\ix n} and \m{e\colon\ix p\sqcup\ix a\to\ix m}
bijections, \m{(\iota_\tau)} a family of finite sets indexed by the infinite places~$\tau$ of~$E$, and $\gamma$ the
weight with \m{\gamma_\tau=|\ix l|-|\iota_\tau|} for all~$\tau$. Let $\mathcal{H}_1$ be the Fock space of the pair of signs
\m{\blockSigns{j}{\cdot}} on~$\ix n$ and of the definite signs on~$\ix a$, and write $v_Q$ also for the Fock vector of a
polynomial~$Q$ in the variables \m{x_{\tau,i,c}}, \m{i\in\ix n}, \m{c\in\ix b}, in the Fock space of the pair of signs
\m{\blockSigns{f}{\cdot}} and \m{\blockSigns{j}{\cdot}}. Then there is such a polynomial~$Q_0$ with the following property.
For every \m{v\in\mathcal{H}_1} there is \m{\varkappa\in\CC} such that for every family \m{\lambda=(\lambda_g)} of
continuous linear forms on~$\mathcal{H}_1$, indexed by \m{g\in G_\infty(\blockSigns{j}{\cdot})}, and every linear
form~$\lambda'$ on the polynomials in these variables with
\begin{enumerate}[label=\textup{(H\arabic*)},leftmargin=2.9em,itemsep=2pt]
\item \m{\lambda_g\bigl(\omega^{\infty}_{\gamma,0}(x)\,w\bigr)=\lambda_{g\,x_1}(w)} for all~$g$, all \m{x=(x_1,x_2)} in
the product of the two groups and all \m{w\in\mathcal{H}_1}, and
\item \m{\lambda_x(v)=\lambda'(Q)} for all~$x$ in \m{G_\infty(\blockSigns{j}{\cdot})} and all polynomials~$Q$ such that
\m{\omega^{\infty}_{\gamma,0}(1,x)\,v_{Q_0}=v_Q},
\end{enumerate}
one has, for every family \m{\eta=(\eta_{\tau k})} of complex \m{\ix p\times\ix m} matrices, \m{k\in\iota_\tau},
\[\varkappa\cdot\lambda_1\bigl(\posFockVector{j}(\sdp{e}{\eta})\bigr)=\lambda'\bigl(\mpoly{f}{j}{\eta}\bigr).\]
\end{proposition}
\leanline{\lean{realmatching}}

In the application \m{\lambda_g(w)} is a theta integral over the definite group with archimedean vector~$w$,
evaluated at~$g$, and \m{\lambda'(Q)} is an integral of a theta kernel with polynomial~$Q$ over the smaller group.
Hypothesis (H2) is then the exchange identity for one vector~$v$ about which nothing explicit is known, and the
conclusion is the same identity for explicit polynomials, at the price of the constant~$\varkappa$.

\begin{proof}
Write \m{d_\tau=|\iota_\tau|}, let \m{\eta^0} be the family with \m{\eta^0_{\tau k}} the \m{\ix p\times\ix m} matrix whose
row~$i$ is the standard basis vector of index \m{e(i)}, for all $\tau$ and~$k$, and put
\m{Q_0=\mpoly{f}{j}{\eta^0}} and \m{v^0=\posFockVector{j}(\sdp{e}{\eta^0})}. Let \m{r\colon\ix a\to\ix n} be the
embedding \m{i\mapsto j(e(i))}. Developing along the rows of~$\eta^0$, the stacked determinant of
\m{\eta^0_{\tau k}} is, up to a sign, the determinant of the \m{\ix a\times\ix a} block of~$x_\tau$ in the columns
\m{e(\ix a)}; so \m{v^0=\pm v_{q_0}} with \m{q_0=\prod_\tau\det(x^{r}_\tau)^{d_\tau}}, where \m{x^{r}_\tau} is the block of
the variables in the columns~$r$.

\emph{Step 1: the constant.} Let $\mathcal{L}$ be the set of the continuous linear forms~$\lambda$
on~$\mathcal{H}_1$ with \m{\lambda(\omega^{\infty}_{\gamma,0}(1,u)w)=\lambda(w)} for all $u$ in the definite group
on~$\ix a$ and all~$w$. For a complex number~$z$ of absolute value~$1$ let \m{t(z)} be the element of the definite
group on~$\ix m$ that is, at every place, diagonal with entries $z^{-2}$ on \m{e(\ix p)} and $z^{-1}$ on
\m{e(\ix a)}, let \m{o(z)} be the scalar~$z^{-1}$ in the definite group on~$\ix a$, and for \m{y\in\mathcal{H}_1} put
\[\Omega(y,z)=\prod_\tau\det(o(z)_\tau)^{d_\tau}\cdot\omega^{\infty}_{\gamma,0}\bigl(\posBlockEmb{j}{t(z)},o(z)\bigr)\,y .\]
\emph{Claim: for every \m{y\in\mathcal{H}_1} there is \m{\varkappa(y)\in\CC} with
\m{\frac{1}{2\pi}\int_0^{2\pi}\lambda\bigl(\Omega(y,e^{i\vartheta})\bigr)\,d\vartheta=\varkappa(y)\,\lambda(v^0)} for all
\m{\lambda\in\mathcal{L}}.} The left side, as a function of~$\lambda$, depends continuously on~$y$ for the
topology of pointwise convergence, because \m{k\mapsto\posBlockEmb{j}{k}} is continuous, and the multiples of the function
\m{\lambda\mapsto\lambda(v^0)} form a closed set of functions on~$\mathcal{L}$. So the set of the~$y$ for which
the claim holds is closed, and since the Fock vectors of polynomials are dense it suffices to prove the claim for
\m{y=v_P}. By Input~6, as in the proof of Lemma~\ref{l:realweil}, the element \m{(\posBlockEmb{j}{t(z)},o(z))} multiplies
\m{v_P} by \m{\prod_\tau\det(o(z)_\tau)^{-d_\tau}}, which cancels the factor in front, and multiplies the variable \m{x_{\tau,i,c}} by \m{z^{-1}\cdot z^{2}=z} if
the column~$c$ lies in \m{j(e(\ix p))}, by \m{z^{-1}\cdot z=1} if it lies in \m{r(\ix a)}, and by~$z$ if it is a negative
column. So \m{\Omega(v_P,z)=v_{P_z}}, where $P_z$ is
$P$ with every variable of a column outside \m{r(\ix a)} multiplied by~$z$. The function
\m{T(z)=\lambda(v_{P_z})} is a polynomial in~$z$, and its mean over the unit circle is \m{T(0)=\lambda(v_q)}, where $q$
is $P$ with the variables outside the columns~$r$ put to zero. For a family~$u$ of unitary matrices, Input~6 gives
\m{\omega^{\infty}_{\gamma,0}(1,u)v_q=\prod_\tau(\det u_\tau)^{-d_\tau}v_{q(u\transp x^{r})}}, so the invariance of~$\lambda$ gives
\m{\lambda(v_{q(u\transp x^{r})})=\lambda(v_q)\prod_\tau(\det u_\tau)^{d_\tau}}. This is the hypothesis of Input~7 for
the index set \m{\ix n=\ix a}, the identity map~$r$, the linear form \m{F\mapsto\lambda(v_F)}, the matrices
\m{Y_\tau=x^{r}_\tau}, the polynomial~$q$ and
\m{K=\lambda(v_q)\prod_\tau\det(X_\tau)^{d_\tau}}. Its conclusion, for the family of identity matrices, is
\m{\ip{q}{q_0}\cdot\lambda(v_{q_0})=\lambda(v_q)\cdot\ip{q_0}{q_0}}. Since \m{\ip{q_0}{q_0}>0} by Input~5\,(c) and
\m{v_{q_0}=\pm v^0}, the claim holds for \m{y=v_P} with \m{\varkappa(v_P)=\pm\ip{q}{q_0}/\ip{q_0}{q_0}}, which does not depend
on~$\lambda$.

\emph{Step 2: the identity for the families \m{\eta^0k\transp}.} Fix~$v$, put \m{\varkappa=\varkappa(v)}, and let $\lambda$,
$\lambda'$ satisfy (H1) and~(H2). Let $k$ be in the definite group on~$\ix m$. By (H1) with \m{x=(1,u)}, the form
\m{\lambda_{\posBlockEmb{j}{k}}} lies in~$\mathcal{L}$. By the claim and by (H1),
\[\varkappa\cdot\lambda_{\posBlockEmb{j}{k}}(v^0)=\frac{1}{2\pi}\int_0^{2\pi}z^{-|\ix a|\sum_\tau d_\tau}\,\lambda_{\posBlockEmb{j}{k\,t(z)}}(v)\,d\vartheta,\qquad z=e^{i\vartheta}.\]
By Input~6 for the pair of signs \m{\blockSigns{f}{\cdot}} and \m{\blockSigns{j}{\cdot}}, the element
\m{(1,\posBlockEmb{j}{k\,t(z)})} sends \m{v_{Q_0}} to \m{v_{Q}} with
\m{Q=\prod_\tau\det(k_\tau t(z)_\tau)^{-d_\tau}\cdot\mpoly{f}{j}{(\eta^0_{\tau k'}\,t(z)_\tau\transp k_\tau\transp)_{\tau,k'}}}: the block
\m{y^{--}} is unchanged and \m{y^{++}} is multiplied on the left by \m{(k\,t(z))\transp}. Since
\m{\eta^0_{\tau k'}t(z)_\tau=z^{-2}\eta^0_{\tau k'}} and \m{\det t(z)_\tau=z^{-2|\ix p|-|\ix a|}}, this is
\m{Q=z^{|\ix a|\sum_\tau d_\tau}\prod_\tau(\det k_\tau)^{-d_\tau}\cdot\mpoly{f}{j}{\eta^0k\transp}}. By (H2) the integrand is
constant, and
\[\varkappa\cdot\lambda_{\posBlockEmb{j}{k}}(v^0)=\prod_\tau(\det k_\tau)^{-d_\tau}\cdot\lambda'\bigl(\mpoly{f}{j}{\eta^0k\transp}\bigr).\]
On the other hand Lemma~\ref{l:realweil}, for $k^{-1}$ and \m{u=1}, and (H1) give
\[\lambda_1\bigl(\posFockVector{j}(\sdp{e}{\eta^0k\transp})\bigr)=\prod_\tau(\det k_\tau)^{d_\tau}\cdot\lambda_{\posBlockEmb{j}{k}}(v^0).\]
Together: the identity of the proposition holds for \m{\eta=\eta^0k\transp}, for every unitary~$k$.

\emph{Step 3: all~$\eta$.} Both sides are values of linear forms at products of determinants \m{\det(\eta_{\tau k}Y_\tau)}.
For the right side this is the definition of the matched polynomial, with \m{Y_\tau=y^{++}_\tau} and the fixed
factor \m{\prod_\tau\det(y^{--}_\tau)^{d_\tau+|\ix n|-2|\ix l|}}. For the left side, a polynomial is determined by its
values, and at every point~$\zeta$ of the space of variables Lemma~\ref{l:detcols} gives matrices
\m{Z_\tau(\zeta)} and numbers \m{c_\tau(\zeta)} with
\m{\stackedDet{\tau}{\eta_{\tau k}}{e}(\zeta)=c_\tau(\zeta)\det(\eta_{\tau k}Z_\tau(\zeta)\transp)} for all~$\eta$. So the
difference of the two sides is \m{M\bigl(\prod_{\tau,k}\det(\eta_{\tau k}Y_\tau)\bigr)} for one linear form~$M$ on the
algebra of pairs of a function of~$\zeta$ and a polynomial, with \m{Y_\tau} the pair \m{(Z_\tau\transp,y^{++}_\tau)}. By
Step~2, $M$ vanishes on these products for the families \m{\eta^0u\transp}, $u$ unitary, which are the values
\m{q_1(u\transp Y)} of the polynomial \m{q_1=\prod_\tau\det(X^{e}_\tau)^{d_\tau}}, with \m{X^{e}_\tau} the rows \m{e(\ix p)} of
a generic \m{\ix m\times\ix p} matrix. Input~7, for the index sets $\ix m$ and~$\ix p$ in the places of $\ix n$
and~$\ix a$, the map~$e$ on~$\ix p$ and \m{K=0}, gives
\m{\ip{q_1}{q_1}\cdot M\bigl(\prod_{\tau,k}\det(\eta_{\tau k}Y_\tau)\bigr)=0} for all~$\eta$, and \m{\ip{q_1}{q_1}>0} by
Input~5\,(c).
\end{proof}

\subsection*{Kernels}
\begin{lemma}[rapid decay of the kernel]\label{l:rapid}
In Setting~\ref{set:bg}, let $S$ and~$V$ be hermitian spaces with finite index sets \m{\ix n} and~\m{\ix b}, $a$ and
$\beta$ characters, $\mathfrak{m}$ a modulus with at least one prime divisor, $\Psi$ a function on the
\m{\ix b\times\ix n} matrices over \m{\OEr/\mathfrak{m}} such that every~$y$ with \m{\Psi(y)\ne0} has a left inverse,
$P$ a polynomial and \m{h\in\UU(V)(\AAA)}. Then \m{g\mapsto\tk{S}{V}{a}{\beta}{\mathfrak{m}}{\Psi}{P}(g,h)} decreases
rapidly on \m{\UU(S)(\AAA)}: its product with every power of the height is bounded.
\end{lemma}
\begin{proof}
The decay estimate of the library bounds the kernel, times any power of the height of~$g$, by a constant times a
moment of the archimedean vector and a power of the height of~$h$, under exactly these hypotheses on
$\mathfrak{m}$ and~$\Psi$; and all moments of a Fock vector~$v_P$ are finite.
\end{proof}
\leanline{\lean{isRapidModKernel}}

\begin{lemma}[from the stabiliser to the full level]\label{l:stab}
In Setting~\ref{set:bg}, let \m{\ix n}, \m{\ix b}, \m{\ix r} be finite index sets, \m{\ix c} an index set, $R'$ a
commutative ring, $\mathfrak{m}$ a modulus, $W$ a hermitian space with index set \m{\ix n\sqcup\ix b}, $\alpha$ and
$\beta$ characters, $\gamma$ a weight and $\chi$ a character of \m{\OEr/\mathfrak{m}} such that $\alpha$ is tame for
\m{(\mathfrak{m},\gamma,\chi)}. Let $P$ be a polynomial with
\m{\omega^{\infty}_{\gamma,0}(1,u)v_P=v_P} for all~$u$ in the definite group on~$\ix b$, let
\m{\varphi\colon\OEr/\mathfrak{m}\to R'} be a ring homomorphism, $w$ a \m{\ix c\times\ix b} matrix over~$R'$,
\m{e\colon\ix r\hookrightarrow\ix b} an embedding, and $F$ a complex function on the \m{\ix b\times\ix n} matrices over
\m{\OEr/\mathfrak{m}} with \m{F(Ay)=F(y)} for all \m{A\in\rowStabilizer{\varphi}{w}{e}} and all~$y$. Write a
\m{\ix b\times(\ix n\sqcup\ix b)} matrix as \m{x=(x'\mid x'')}, and let $\Gamma$ be the level of the elements of
\m{\KK{\ix b}} whose reduction modulo~$\mathfrak{m}$ lies in \m{\rowStabilizer{\varphi}{w}{e}}. Then for every
\m{g\in\UU(W)(\AAA)}
\[\clsum{h}{\Gamma}\tk{W}{I_{\ix b}}{\alpha}{\beta}{\mathfrak{m}}{\Psi_1}{P}(g,h)=\clsum{h}{\KK{\ix b}}\tk{W}{I_{\ix b}}{\alpha}{\beta}{\mathfrak{m}}{\Psi_0}{P}(g,h),\]
with \m{\Psi_1(x)=\columnDatum{\chi}{\varphi}{w}{e}{x''I_e}\cdot[\FixesColsOff{e}{x''}]\,F(x')} and
\m{\Psi_0(x)=\chi(\det x'')\,F(x''^{-1}x')}.
\end{lemma}
\begin{proof}
The kernel is invariant under left translation of~$h$ by rational points: the theta functional is (Input~8\,(b)),
and the archimedean component of a rational point fixes~$v_P$ by the hypothesis on~$P$. By Input~8\,(h) the left
side is the class sum at level \m{\KK{\ix b}} of the sum over the cosets~$q$ of the stabiliser in
\m{\GL_{\ix b}(\OEr/\mathfrak{m})} of the translates by lifts of~$q$. By Input~8\,(d) and the tameness of~$\alpha$,
the translate by a lift of~$q$ is the kernel with datum \m{x\mapsto\chi(\det q)\Psi_1(q^{-1}x)}. By
Lemma~\ref{l:cols}, with \m{z=q^{-1}x''} and \m{y=q^{-1}x'},
\[\chi(\det q)\,\Psi_1(q^{-1}x)=\bigl[q^{-1}x''\in\rowStabilizer{\varphi}{w}{e}\bigr]\,\chi(\det x'')\,F(x''^{-1}x').\]
If $x''$ is invertible, exactly one coset~$q$ contains~$x''$; if not, \m{\chi(\det x'')=0}. So the sum over~$q$ of
the data is~$\Psi_0$, and the kernel is linear in its datum.
\end{proof}
\leanline{\lean{classSumRowStabilizer}}

\begin{proposition}[matching summed against a form on a square block]\label{p:sum}
In Setting~\ref{set:bg}, let \m{\ix n} and~\m{\ix b} be finite index sets, $V$ a hermitian space with index
set~$\ix b$ and $U$ one with index set~$\ix n$, $\mathfrak{m}$ a modulus, $a$, $\beta$, $\beta'$ characters, $e$ a
weight and $\chi$ a character of \m{\OEr/\mathfrak{m}} such that \m{\beta\beta'} is tame for \m{(\mathfrak{m},e,\chi)}.
Let $\Xi$ be a function on \m{\UU(I_{\ix b})(\AAA_f)}, $F$ a function on \m{\UU(V)(\AAA)} and \m{\varkappa\in\CC} with
\begin{equation}\label{e:hF}
\varkappa\cdot\clsum{h}{\KK{\ix b}}\Xi(h)\,\tk{V}{I_{\ix b}}{a}{\beta}{\mathfrak{m}}{\chi\circ\det}{1}(x,h)=F(x)\qquad\text{for all }x\in\UU(V)(\AAA).
\end{equation}
Let $\mathfrak{n}$ be a modulus with \m{\mathfrak{m}\subseteq\mathfrak{n}}, $f$ a function on the \m{\ix n\times\ix b}
matrices over \m{\OEr/\mathfrak{n}}, $P$ a polynomial, \m{g\in\UU(U)(\AAA)}, \m{\varkappa'\in\CC}, and let $K$ assign to
a function~$\dL$ and a point \m{h\in\UU(I_{\ix b})(\AAA_f)} a function \m{K(\dL,h)} on \m{\UU(I_{\ix n})(\AAA_f)}
such that, for all partners $\dL$, $\Psi$ for the identity of \m{\ix b\sqcup\ix n} and~$\chi^{-1}$ and all~$h$,
\begin{equation}\label{e:hK}
\varkappa'\cdot\clsum{u}{\KK{\ix n}}K(\dL,h)(u)=\int_{[\UU(V)]}\tk{V}{I_{\ix b}\oplus U}{a}{\beta\beta'}{\mathfrak{m}}{\Psi}{P}\bigl(x,(h,g)\bigr)\,dx,
\end{equation}
the polynomial~$P$ being read in the variables of the rows of~$U$. Then, with
\[\dL_0(x'\mid x'')=\chi(\det x'')\sum_{k\in\GL_{\ix b}(\OEr/\mathfrak{n})}f\bigl((x''^{-1}x'\bmod\mathfrak{n})\,k\bigr),\]
one has
\begin{multline*}
\varkappa'\cdot\frac{\varkappa}{|\GL_{\ix b}(\OEr/\mathfrak{n})|}\cdot\clsum{u}{\KK{\ix n}}\ \clsum{h}{\KK{\ix b}}\Xi(h)\,K(\dL_0,h)(u)\\
=\int_{[\UU(V)]}F(x)\cdot\tk{V}{U}{a}{\beta'}{\mathfrak{n}}{f}{P}(x,g)\,dx .
\end{multline*}
\end{proposition}
\begin{proof}[Proof \textup{(}the translation rule is described, not recomputed\textup{)}]
The proof is complete but for the exact form of the translation rule of Input~8\,(d), which the formal text applies
by simplification. For a matrix~$y$ with square upper block~$y_1$ and lower block~$y_2$ let
\m{\Psi(y)=\chi(\det y_1)\,f(y_2\bmod\mathfrak{n})}, and let
\[\Psi_0=\sum_{k}\chi(\det\tilde k)^{-1}\,\Psi(\,\cdot\,\tilde k),\]
the sum over \m{k\in\GL_{\ix b}(\OEr/\mathfrak{n})} with a lift~$\tilde k$ of~$k$ to \m{\GL_{\ix b}(\OEr/\mathfrak{m})}. By
Input~9\,(b), applied to the function \m{\sum_kf(\,\cdot\,k)} and after a change of the summation index, $\dL_0$
and~$\Psi_0$ are partners for~$\chi^{-1}$. So \eqref{e:hK} holds for them. A matrix~$y$ with \m{\Psi(y)\ne0} has
$y_1$ invertible, hence a left inverse, so by Lemma~\ref{l:rapid} the kernels with datum~$\Psi$ are integrable
over \m{[\UU(V)]}. By Input~8\,(d) the kernel with datum \m{\chi(\det\tilde k)^{-1}\Psi(\,\cdot\,\tilde k)} is the kernel
with datum~$\Psi$ at the point~$x$ translated by an integral point, and the integral over \m{[\UU(V)]} does not
change under such a translation. So the right side of \eqref{e:hK} is \m{|\GL_{\ix b}(\OEr/\mathfrak{n})|} times the
integral of the kernel with datum~$\Psi$, which by the seesaw (Input~8\,(e)) and the change of modulus
(Input~8\,(f)) is
\[\tk{V}{I_{\ix b}}{a}{\beta}{\mathfrak{m}}{\chi\circ\det}{1}(x,h)\cdot\tk{V}{U}{a}{\beta'}{\mathfrak{n}}{f}{P}(x,g).\]
Multiply by \m{\varkappa\,\Xi(h)}, take the class sum over~$h$, which is a finite sum and commutes with the integral,
and use~\eqref{e:hF}.
\end{proof}
\leanline{\lean{explicitMatchingOfSum}}

\subsection*{Proofs of the theorems}
\begin{proof}[Proof of Theorem A]
The signs of~$W$ are \m{\blockSigns{j}{\cdot}}, so $W$ has \m{|\ix l|} negative signs, as has~$V'$, and \m{I_{\ix a}} has
none. Take $Q_0$ and the function~$\varkappa$ of Proposition~\ref{p:real} for $f$, $j$, $e$, \m{(\iota_\tau)}. Apply
Input~9\,(a) to $W$, \m{V=I_{\ix a}}, $V'$, $\sigma$ and the vector \m{v=v_{Q_0}}, whose moments are finite; its
numerical hypotheses hold by the assumptions on \m{|\ix a|}, \m{|\ix l|}, \m{|\ix n|} and by the weight of the
character below. It gives $C_0$ and~$x_0$. Put \m{C=|\varkappa(x_0)|^2/C_0}.

Let $\mathfrak{m}$, $\alpha$, $\gamma$, $\chi$ be as in the statement. The character
\m{\nu=\alpha^{-1}\upsilon^{|\ix b|+d}} is tame for $\mathfrak{m}$ and~$\chi^{-1}$, and
\m{\nu^{-1}\upsilon^{\lfloor|\ix n|/2\rfloor}=\alpha}, because \m{|\ix n|=|\ix a|+|\ix b|=2(|\ix b|+d)}. Input~9\,(a)
gives \m{c_0\ne0} with \m{|c_0|^2=C_0\,\LP{\mathfrak{m}}{\nu}{d}}; put \m{c=\varkappa(x_0)/c_0}, so that $|c|^2$ is as
stated. Fix~$\eta$, partners $\dL$, $\Psi$ for \m{(\sigma,\chi^{-1})}, $\beta$ and~$g$. For \m{g'\in G_\infty(\blockSigns{j}{\cdot})},
regarded as an adelic point with trivial finite part, and a vector~$w$ of the Fock space, let
\begin{align*}
\lambda_{g'}(w)&=\int_{[\UU(I_{\ix a})]}\theta\bigl(\omega_{\alpha,1}(gg',h)(w\otimes\dL)\bigr)\,dh,\\
\lambda'(Q)&=c_0\int_{[\UU(V')]}\tk{V'}{W}{\alpha}{\alpha}{\mathfrak{m}}{\Psi}{Q}(x,g)\,dx .
\end{align*}
The forms~$\lambda_{g'}$ are continuous, and \m{\lambda'} is defined because the kernels are rapidly decreasing
(Lemma~\ref{l:rapid}; a matrix in the support of~$\Psi$ consists of columns of an invertible matrix). Hypothesis
(H1) holds because the theta integral intertwines the Weil representation at the infinite places with right
translation. Hypothesis (H2), for \m{v=x_0}: if \m{\omega^{\infty}_{\gamma,0}(1,g')v_{Q_0}=v_Q}, then by the exchange
identity at the point~$gg'$
\begin{align*}
\lambda_{g'}(x_0)&=c_0\int_{[\UU(V')]}\theta\bigl(\omega_{\alpha,\alpha}(x,gg')(v_{Q_0}\otimes\Psi)\bigr)dx\\
&=c_0\int_{[\UU(V')]}\theta\bigl(\omega_{\alpha,\alpha}(x,g)(v_{Q}\otimes\Psi)\bigr)dx=\lambda'(Q).
\end{align*}
Proposition~\ref{p:real} gives
\m{\varkappa(x_0)\,\lambda_1\bigl(\posFockVector{j}(\sdp{e}{\eta})\bigr)=\lambda'\bigl(\mpoly{f}{j}{\eta}\bigr)}. The vector
\m{\posFockVector{j}(\sdp{e}{\eta})} is fixed by the definite group at infinity (Lemma~\ref{l:realweil} with \m{k=1}), and
the kernel is invariant under right translation of~$h$ by \m{\KK{\ix a}}: by Input~8\,(d) the translation acts on
the datum through~$\chi$, and the partner~$\dL$ absorbs it (Input~9\,(b)). So by Input~8\,(h) the integral
\m{\lambda_1} is the class sum, and after division by~$c_0$
\[c\cdot\clsum{h}{\KK{\ix a}}\tk{W}{I_{\ix a}}{\alpha}{1}{\mathfrak{m}}{\dL}{\sdp{e}{\eta}^{j}}(g,h)
=\int_{[\UU(V')]}\tk{V'}{W}{\alpha}{\alpha}{\mathfrak{m}}{\Psi}{\mpoly{f}{j}{\eta}}(x,g)\,dx .\]
Twisting by~$\beta$ in the second place on the left and by~$\beta^{-1}$ in the first place on the right multiplies
both sides by \m{\beta(\det g)^{-1}} (Input~8\,(c)), which gives the statement.
\end{proof}

\begin{proof}[Proof of Theorem B]
Apply Theorem~A with \m{\ix n=\ix m=\ix a=\{1,\dots,m\}}, \m{\ix p=\ix l=\ix b=\emptyset}, \m{W=I_m}, $V'$ the zero
space, \m{d=m/2} and \m{\iota_\tau} a set of $c_\tau$ elements. The conditions \m{|\ix a|=2d} and \m{2\le m} hold since
$m$ is even and \m{m\ge1}, and \m{\gamma_\tau=-c_\tau}. Take the family~$\eta$ in which each of the $c_\tau$ matrices \m{\eta_{\tau k}}, \m{k\in\iota_\tau}, is the empty
\m{0\times m} matrix: a matrix with no rows stacked
over~$x_\tau$ is~$x_\tau$, so \m{\sdp{e}{\eta}=\prod_\tau\det(x_\tau)^{c_\tau}}, and the matched polynomial is~$1$. Take
the partners of Input~9\,(b) for \m{f=1}: \m{\dL=\chi\circ\det} and \m{\Psi=1}. Take \m{\beta=\alpha}. The right side of
Theorem~A is the integral of the constant~$1$ over the quotient of the trivial group, a positive number~$\varrho_0$.
So Theorem~B holds with \m{\gamma=c/\varrho_0} and the constant \m{C/\varrho_0^2}.
\end{proof}

\begin{proof}[Proof of Theorem C \textup{(}the characters are not recomputed\textup{)}]
Three facts are combined; the characters in the seesaw and in the negation of the first space are followed as the
formal text follows them, by simplification, and were not recomputed. Write \m{J=I_m} and \m{S_0=I_{2m,m}}. First,
by Input~8\,(g) there is~$n_0$ with
\m{|n_0|=1} relating the kernels of \m{(J^{-},J)} to the conjugates of those of \m{(J,J)}. Second, Theorem~B with
\m{c_\tau=0} gives a constant~$C_1$ and, for the character~$c^{-1}$, which is tame of weight~$0$ for
\m{(\mathfrak{m},\chi^{-1})}, a number~$\gamma_1$ with
\m{|\gamma_1|^2=C_1/\LP{\mathfrak{m}}{c\,\upsilon^{m/2}}{m/2}} and
\[\gamma_1\clsum{h}{\KK{m}}\tk{J}{J}{c^{-1}}{c^{-1}}{\mathfrak{m}}{\chi^{-1}\circ\det}{1}(g,h)=1 .\]
Conjugating and passing
to~$J^{-}$ gives \eqref{e:hF} for \m{V=J^{-}}, the functions \m{\Xi=1} and \m{F=1}, and
\m{\varkappa=\bar\gamma_1n_0/\chi(\epsilon)^m}, where $\epsilon$ is the element of norm~$-1$ of Input~8\,(g) and
\m{|\chi(\epsilon)|=1}. Third, Theorem~A for \m{W=J\oplus S_0}, with \m{4m} indices of which $m$ are negative,
\m{V'=J^{-}}, \m{|\ix a|=3m}, \m{d=m} and empty sets~$\iota_\tau$, so that \m{\gamma_\tau=m}, gives~$C_2$ and, for the
character \m{c\,\upsilon^{m}}, tame of weight~$m$, a number~$\gamma_2$ with
\m{|\gamma_2|^2=C_2/\LP{\mathfrak{m}}{c^{-1}\upsilon^{m}}{m}}; its matched polynomial is
\m{\prod_\tau\det(y^{--}_\tau)^{2m}}, the polynomial of the carrier, and its identity is \eqref{e:hK} with
\m{\varkappa'=\gamma_2}, \m{U=S_0} and \m{K(\dL,h)(u)} the kernel of \m{(J\oplus S_0,I_{3m})} at \m{((h,x),u)}.
Proposition~\ref{p:sum}, with the modulus \m{\mathfrak{n}=\unitModulus{E}}, whose general linear group has one
element, and \m{f=1}, gives
\begin{align*}
\gamma_2\,\varkappa\cdot\clsum{u}{\KK{3m}}\ \clsum{h}{\KK{m}}K(\dL_0,h)(u)
&=\int_{[\UU(J^{-})]}\tk{J^{-}}{S_0}{\upsilon^{m}}{1}{\unitModulus{E}}{1}{\textstyle\prod_\tau\det(y^{-}_\tau)^{2m}}(x',x)\,dx'\\
&=\carrier{E}{m}(x),
\end{align*}
with \m{\dL_0(x'\mid x'')=\chi(\det x'')}. By the seesaw in the first place (Input~8\,(e)) the kernel
\m{K(\dL_0,h)(u)} is the product of the kernel of \m{(J,I_{3m})} with constant datum at \m{(h,u)} and of the kernel
of \m{(S_0,I_{3m})} with datum \m{\chi\circ\det} at \m{(x,u)}; the class sum of the first over~$h$ is
\m{\carrierXi{m}{c}{u}}. This is Theorem~C with \m{\gamma=\gamma_2\varkappa} and \m{C=C_1C_2}.
\end{proof}
\leanline{\lean{explicitmatching}\sep \lean{classSumDet}\sep \lean{carrierIdentity}}

\sect{Where the results are used}{s:used}
The auxiliary pairs, their forms and their descent are used throughout Chapters VII and~VIII: the single lifts of
the chain are adjoint to descents of auxiliary pairs. The divided splitting, the stabiliser and the column datum
serve the glued lattices of the note \emph{Glued lattices, product forms and the integral comparison forms}.
Theorem~A, Lemma~\ref{l:stab}, Proposition~\ref{p:sum} and Theorem~C are the inputs of the identity ``the twist
survives only as a constant'', which is Theorem~C of the note \emph{The height of a ratio of lift volumes};
Lemma~\ref{l:coord} serves the line multipliers of that note, and Lemma~\ref{l:wnprod} gives there the size of the
polynomials.

\sect{Notes on the formal text}{s:notes}
\begin{enumerate}[label=\textup{\arabic*.},leftmargin=1.6em,itemsep=5pt]
\item \emph{Arguments that only a compiler can settle.} The arguments of the declarations were read from the text.
For six definitions and two results, a lemma and a proposition, this reading lists a letter that the definition binds itself, as the index
of a product or as the variable of a quantifier, and that Lean, by its rule for section variables, presumably
does not take as an argument: the place~$\tau$ in the product of stacked determinants, in the substitution
\m{\sub{c}{g}}, in \m{\infc}, in the matched polynomial, in the carrier, in Lemma~\ref{l:subst} and in
Proposition~\ref{p:real}; the map~$j$ in \m{\sub{c}{g}}; the family~$\eta$ in Proposition~\ref{p:real}, where it is
quantified in the conclusion; and the condition \m{w\ne\bar w} in the definition of ``lifts with parameters'',
where it is quantified inside. The uses of these declarations in the files confirm the shorter lists (the carrier
is written with the two arguments $E$ and~$m$, the matched polynomial with $f$, $j$, $\eta$). This note states
the declarations with the arguments that their uses show; it was not confirmed by asking Lean. For this reason
these six definitions carry no notation made by the build.
\item \emph{Instances listed twice.} In the lemmas of Remark~\ref{r:book} and in Lemmas \ref{l:red}
and~\ref{l:genres} the text reading lists the assumptions on~$E$ twice, once from a variable line that a later one
shadows.
\item \emph{An argument that is not used.} In Lemma~\ref{l:genres} the divided splitting enters only through the
element \m{\genResidue{\mathcal{V}}}, which depends on~$O$ alone; in Definition~\ref{d:div}\,(iii) its argument is
unnamed in the formal text.
\item \emph{Constants that may vanish.} Theorems A to~C assert the existence of a real number~$C$ without sign or
non-vanishing. If \m{C=0}, the identities say that their right sides vanish. In Theorem~A the formal constant is
\m{|\varkappa(x_0)|^2/C_0}, with $C_0$ from the exchange theorem.
\item \emph{Truncated subtraction.} The exponent \m{|\iota_\tau|+|\ix n|-2|\ix l|} of the matched polynomial, the
rank difference \m{|\ix n|-|\ix r|} and the exponent in the test value are differences of natural numbers, which
Lean truncates at zero; under the hypotheses of the statements they are not negative.
\item \emph{Index sets of Theorem~C.} The second space of the kernel in Theorem~C has the index set of
\m{I_{2m,m}}, the disjoint union of \m{\{1,\dots,2m\}} and \m{\{1,\dots,m\}}; it is written \m{I_{3m}} above.
\item \emph{What was not reconstructed.} In the proof of Theorem~C the bookkeeping of the characters in the seesaw
and in the negation of the first space, and in the proof of Proposition~\ref{p:sum} the exact form of the
translation rule, are carried out in the formal text by simplification; they are described here and were not
recomputed. In Step~3 of the proof of Proposition~\ref{p:real} the linear form~$M$ is defined in the formal text on
pairs, after multiplication by the pair of the product of the numbers \m{c_\tau(\zeta)^{d_\tau}} and of the fixed
factor of the matched polynomial; the proof above says this in words.
\item \emph{Letters that differ from the formal text.} The partner on the large side is a capital phi in the formal
text and~$\dL$ here, as in the notes on the unit \emph{Exchange}; the finite set of places is a capital phi there
and~$\Plc$ here; the level places are~$S$ there and~$\lpl$ here; the second ring of §\ref{s:cols} and of Lemma~\ref{l:stab}
is~$S$ there and~$R'$ here; the matrix of Lemma~\ref{l:genres} is~$S$ there and~$Z$ here; the matrix of the column
datum is~$\sigma$ there and~$x$ here; the constants of Propositions \ref{p:real} and~\ref{p:sum} are~$\kappa$,
$\kappa'$ there and~$\varkappa$, $\varkappa'$ here. The capital phi and the letters $S$, $\sigma$ and~$\kappa$ have other
meanings in this write-up.
\end{enumerate}

\subhead{How this note was written.} From the dossier of the unit, whose reader calls the step claims of the long
proofs a first reading and notes that no second reader had looked at them. After the dossier was frozen, the second reader of the write-up compared the readings of all statements of the unit
with the Lean text and sampled its structured proofs (her message to the writer of 26 November 2026); her two
remarks on this part, the rank of the second space in Theorem~C and the direction of the space of homomorphisms
behind the Fock vectors of Definition~\ref{d:real}, agree with the statements above. I read the five Lean files of these
declarations completely. The proof of the archimedean matching (Proposition~\ref{p:real}) I read line by line in
the Lean text and recomputed its character identities by hand from Input~6: the cancellation of the determinant
factor in~$\Omega$, the exponents of~$z$ in Steps 1 and~2, and the two uses of polarisation. These computations
are mine; the formal proof obtains them by simplification, and I could not see its intermediate goals. I read the
Lean statements of the exchange theorem, of the polarisation lemma, of the rule for the real Weil representation
on Fock vectors, of the notions of a tame character, of partners and of the product of $L$-values, and the
docstrings of the other definitions of Setting~\ref{set:bg}; I did not read the proofs of the inputs. The proofs of
Lemmas \ref{l:cols}, \ref{l:detcols} and~\ref{l:subst} are written out by hand and are shorter than the formal ones.
The program check of the unit's dossier covered \UnitSteps{} leaf steps of \UnitProofs{} proofs.

\subhead{Revision.} This text was revised after the report of an independent referee (26 November 2026), who
compared the statements with the Lean text and recomputed the proofs of Proposition~\ref{p:real} and of
Theorems A to~C. The changes: the partial $L$-function is that of the base change of the character, without the
Euler factors at the conjugate places as well; Input~7 is stated with the sizes of its matrices; the box says which
two computations are described and not carried out; several letters were changed, among them the partner on the
large side, now~$\dL$ as in the notes on the unit \emph{Exchange}.

\head{Appendix A. Concordance with the Lean artifact}
{\fontsize{8.6}{11.6}\selectfont\renewcommand{\leanfont}{\ttfamily\fontsize{7.7}{10}\selectfont\color{leancol}}
\setlength{\LTleft}{0pt}\setlength{\LTright}{\fill}\setlength{\tabcolsep}{0pt}
\begin{longtable}{@{}p{30mm}p{56mm}p{43mm}r@{}}
\textbf{In this note} & \textbf{Lean declaration} & \textbf{File} & \textbf{Lines}\\\midrule\endhead
\ConcordanceRows
\end{longtable}}

\subhead{The three theorems and the archimedean matching as they are written in the artifact.}
\quoted{explicitmatching}
\quoted{classSumDet}
\quoted{carrierIdentity}
\quoted{realmatching}
In these texts {\leanfont modKernel} and {\leanfont modKernelI} are the theta kernels, {\leanfont stdClassSum} the
class sum, {\leanfont Kf} the stabiliser of the standard lattice, {\leanfont IsTameFor} tameness, {\leanfont IsPartner}
the partners, {\leanfont lValueProd} the product of $L$-values, {\leanfont upsilonU1} the character~$\upsilon$,
{\leanfont realPairWeil} the real Weil representation and {\leanfont polyVector} the Fock vector of a polynomial.

\head{Notation}
\begingroup\fontsize{8.6}{10.9}\selectfont\raggedright
\setlength{\columnsep}{9mm}\setlength{\parskip}{0pt}
\begin{multicols}{2}
\nt{$E$, $\OEr$, $\Ohat$, $\AAA$, $\AAA_f$}{a CM field; its integers; integral finite adeles; adeles; finite adeles}
\nt{$\tau$, $w$, $\bar w$, $k_w$, $q_w$}{an infinite place; a finite place, its conjugate, residue field and norm}
\nt{$\upsilon$}{the character of the Weil data}
\nt{$\ix r,\ix n,\ix N,\ix m,\ix a,\ix b,\ix c,\ix p,\ix l,\ix L,\ix k$, $\ix n'$}{finite index sets; \m{|\ix n|}
the number of elements}
\nt{$\mathfrak{m}$, $\mathfrak{n}$}{moduli}
\nt{$\lpl$, $\mathfrak{M}_\lpl$, $\mathfrak{Q}_\lpl$}{level places; the products of the chosen and of the
twist primes}
\nt{$\alpha,\beta,\beta',\mu,\mu',\nu$, $a$, $c$}{characters of the norm-one ideles}
\nt{$\gamma,\gamma'$, $e$, $\mu$ (weights)}{integer weights at the infinite places}
\nt{$\psi$, $\chi$}{a tame part; a multiplicative character of a residue ring}
\nt{$\LP{T}{\nu}{d}$, $L^{T}(z,\nu)$}{the product of $L$-values; the partial Hecke $L$-function of the base change
of~$\nu$, without the Euler factors at~$T$ and at the conjugate places}
\nt{$s$, $s'$}{sign vectors}
\nt{$S$, $T$, $U$, $V$, $V'$, $W$, $J$, $S_0$}{hermitian spaces}
\nt{$I_{\ix n}$, $I_m$, $I_{2m,m}$, $S^{-}$}{the standard definite space; the signed space; the opposite space}
\nt{$\UU(S)(\AAA)$, $[\UU(S)]$, $G_\infty(s)$}{adelic unitary group; automorphic quotient; archimedean group}
\nt{$\KK{\ix n}$, $\Gamma$}{the stabiliser of the standard lattice; a level}
\nt{$\clsum{h}{\Gamma}$, $\mathfrak{c}$, $h_{\mathfrak{c}}$}{the class sum; a class; its representative}
\nt{$\omega_{\alpha,\beta}$, $\omega^{\infty}_{\gamma,\gamma'}$, $\theta$}{the Weil representation; its archimedean part; the
theta functional}
\nt{$v_P$, $x_{\tau,i,c}$}{the Fock vector of a polynomial~$P$; a variable}
\nt{$\tk{S}{T}{\alpha}{\beta}{\mathfrak{m}}{\Psi}{P}(g,h)$}{the theta kernel}
\nt{$\dL$, $\Psi$, $\Psi_0$, $\Psi_1$, $\dL_0$, $f$ (datum)}{functions on matrices over a residue ring; $\dL$,
$\Psi$ partners}
\nt{$\Plc$}{a finite set, in the applications the set of infinite places}
\nt{$\sigma$}{a bijection \m{\ix b\sqcup\ix a\to\ix n}}
\nt{$\kappa$, $\kappa'$, $b_l$, $G_\kappa$, $h_\kappa$, $S_\kappa(z)$}{weights; basis polynomials; Gram matrix; hook
number; Schur matrix}
\nt{$\operatorname{pol}_\kappa(\xi)$, $\|\xi\|_\kappa$, $\xi$, $\xi'$, $\xi''$}{the polynomial and the norm of a coordinate
vector}
\nt{$\ip{P}{Q}$}{the Fock pairing}
\nt{$R$, $R'$, $\varphi$, $w$ (matrix), $I_e$, $I_u$}{rings, a homomorphism, a matrix; inclusion matrices}
\nt{$\FixesColsOff{e}{z}$, $\rowStabilizer{\varphi}{w}{e}$, $\columnDatum{\chi}{\varphi}{w}{e}{x}$}{Definition~\ref{d:cols}}
\nt{$\HasWeightNormLe{p}{\kappa}{F}{b}$}{Definition~\ref{d:wn}}
\nt{$\ChosenGL{\lpl}{\ix r}$, $\firstBlockHom{\lpl}{u}{g}$, $\auxLevel{\lpl}{u}$}{Definition~\ref{d:levels}}
\nt{$O$, $\AuxPair{E}{\ix r}{\ix n}$}{an auxiliary pair; the set of them}
\nt{$\mathfrak{M}_O$, $\mathfrak{Q}_O$, $s_O$, $W_O$, $e_O$, $\operatorname{ord}_O$, $\kappa^{\pm}_O$, $\mu_O$,
$\psi_O$}{its components}
\nt{$\levelGroup{O}$, $\levelChar{O}$, $\auxLevelChar{O}{g}$, $\resDetDatum{O}{X}$}{level, characters, datum}
\nt{$\auxKappa{O}$, $\detExp{O}{\tau}$, $\FormV{O}$, $\rankDiff{O}$}{weight, determinant exponent, forms, rank
difference}
\nt{$\chosenGLRep{O}{\pi}$, $\formAction{O}$}{representations of \m{\ChosenGL{O}{\ix r}}}
\nt{$\tensorPoly{O}$, $\harmonicPoly{O}$}{the tensor polynomial; the harmonic polynomial}
\nt{$\auxA{O}$, $\auxB{O}$, $r_-$, $\pairThetaKernel{O}$}{the characters of the kernel; the number of negative
signs; the kernel}
\nt{$\auxDatum{O}$, $\fockDatum{O}{\xi}$, $\conjKernel{O}{g}{h}$, $\descent{O}{\Xi}{g}$}{finite datum, test datum,
conjugate kernel, descent}
\nt{$H_{O,w}(X)$, $T_{w,i}$, $\testValue{O}{w}$, $\nu_\alpha(w)$, $z$ (parameters)}{Hecke test}
\nt{$\chosenModulusSq{\lpl}$, $\reduceToRzero{x}$, $\unitModulus{E}$}{Definition~\ref{d:moduli}}
\nt{$\genIdele{O}$, $\genIntAdele{O}$, $\genResidue{\mathcal{V}}$}{the generator of \m{\mathfrak{M}_O}; its class}
\nt{$\mathcal{V}$, $\DividedSplitting{\ix N}{\ix m}{\ix a}{O}{z}$, $\varepsilon$}{a divided splitting; the set of them; its
bijection}
\nt{$\stackedDet{\tau}{\eta}{j}$, $\sdp{e}{M}$, $\sdp{e}{\eta}^{j}$}{stacked determinant; the product; with
renamed columns}
\nt{$\sub{c}{g}$, $\infc$, $\firstColsCoord{\kappa}{F}{l}$}{substitution; first columns; coordinates}
\nt{$\eta$, $\eta^0$, $M$, $\iota_\tau$, $d_\tau$}{families of matrices; the index sets and their sizes}
\nt{$\carrier{E}{m}$, $\carrierXi{m}{c}{h}$, $y^{-}_\tau$}{the carrier; the class sum of Definition~\ref{d:carrier}}
\nt{$\blockSigns{f}{i}$, $\posBlockEmb{j}{k}$, $\posFockVector{j}$}{block signs; the positive block; Fock vectors}
\nt{$\mpoly{f}{j}{\eta}$, $y^{++}_\tau$, $y^{--}_\tau$}{the matched polynomial; blocks of variables}
\nt{$f$, $j$, $e$ (bijections)}{bijections of index sets}
\nt{$d$, $m$, $p$, $h$, $i$, $k$, $l$}{integers, indices, fillings; $h$, $g$, $x$, $u$ also points of groups}
\nt{$C$, $C_0$, $C_1$, $C_2$, $c$, $c_0$, $\gamma$, $\gamma_1$, $\gamma_2$, $\varkappa$, $\varkappa'$, $n_0$,
$\varrho_0$}{constants}
\nt{$\mathcal{H}_1$, $\mathcal{L}$, $\lambda$, $\lambda_g$, $\lambda'$, $Q_0$, $Q$, $v$, $v^0$, $x_0$}{Fock space,
invariant forms, linear forms, polynomials, vectors of Proposition~\ref{p:real}}
\nt{$t(z)$, $o(z)$, $\Omega(y,z)$, $T(z)$, $P_z$, $q$, $q_0$, $q_1$, $r$, $\vartheta$}{objects of its proof;
$\vartheta$ an angle}
\nt{$X$, $X'$, $Y$, $Z$, $U_1$, $U_2$, $A$, $B$, $K$, $F$, $F'$, $\Xi$, $\zeta$, $\epsilon$}{matrices, polynomials,
functions, points}
\nt{$\mathcal{S}$, $M$ (linear form), $\pi$, $y$, $a$, $b$, $u$, $g$, $x$, $x'$, $x''$, $y_1$, $y_2$, $\tilde k$}{an
algebra; a linear form; a representation; vectors, matrices, blocks}
\nt{$L$}{in ``$L$-function'' and ``$L$-value''}
\end{multicols}
\endgroup

\end{document}
